% Traduction : « avoir beau » ; « demander » ; « depuis » — Espagnol Tle A — Cours signé OnBuch+
% Programme officiel MINESEC. Compiler avec : tectonic E25-aunque-pedir-desde-hace.tex
\newcommand{\DOCMATIERE}{Espagnol}
\newcommand{\DOCNIVEAU}{Tle A}
\newcommand{\DOCMODULE}{Módulo 5 : Médias, communication et culture de la paix}
\newcommand{\DOCLECON}{E25}
\newcommand{\DOCTITRE}{Traduction : « avoir beau » ; « demander » ; « depuis »}
% OnBuch+ — préambule « cours signés » ENRICHI (compilé avec tectonic / XeLaTeX)
% Système visuel « Pop » d'OnBuch+ : fond crème (jamais blanc), bordures encre
% épaisses, ombres « dures » décalées, titres Archivo Black en capitales.
% Version MATHÉMATIQUES : graphes (pgfplots/tikz), boîtes pédagogiques
% (définition, propriété, méthode, attention, le savais-tu, exercice résolu,
% exercice, TP, bilan), page de garde et signature OnBuch+ ancrée en bas.
%
% Macros attendues dans main.tex AVANT \input{preamble} :
%   \DOCMATIERE  \DOCNIVEAU  \DOCMODULE  \DOCLECON (numéro)  \DOCTITRE
\documentclass[11pt]{article}

\usepackage{fontspec}
\usepackage[spanish,french]{babel} % français = langue principale ; l'espagnol via \esp{..} / espagnol
\usepackage[a4paper,margin=2cm,top=3.4cm,bottom=2.8cm,headheight=1.4cm,footskip=1.3cm]{geometry}
\usepackage{amsmath,amssymb,amsfonts}
\usepackage{mathtools} % \xmapsto, \coloneqq... souvent utilisés par les modèles
\usepackage{cancel} % simplifications barrées (bilans de forces)
% \abs{z} (module, valeur absolue) et \norm{u} (norme), utilisés par les modèles
\providecommand{\abs}{}\let\abs\relax\DeclarePairedDelimiter{\abs}{\lvert}{\rvert}
\providecommand{\norm}{}\let\norm\relax\DeclarePairedDelimiter{\norm}{\lVert}{\rVert}
\usepackage[table]{xcolor} % [table] : fournit aussi \rowcolors (alternance de lignes)
\usepackage{pagecolor}
\usepackage{tikz}
\usetikzlibrary{arrows.meta,positioning,calc,shapes.geometric,shapes.misc,shapes.arrows,decorations.pathmorphing,decorations.pathreplacing,decorations.markings,patterns,shadows,mindmap,trees,fit,backgrounds,matrix,intersections,angles,quotes}
\usepackage{pgfplots}
\usepackage[version=4]{mhchem} % équations nucléaires \ce{^{235}_{92}U}
\usepackage[european,siunitx]{circuitikz} % schémas électriques
\pgfplotsset{compat=1.18}
\usepgfplotslibrary{fillbetween,groupplots} % aires entre courbes (calcul intégral)
\usepackage{enumitem}
\usepackage{fancyhdr}
% [most] charge toutes les bibliothèques tcolorbox (listings, minted, raster,
% documentation...) et gonfle inutilement la mémoire TeX sur un document long
% (~300 boîtes) : on ne charge que ce qui est réellement utilisé.
\usepackage[skins,breakable]{tcolorbox}
\usepackage{graphicx}
\usepackage{colortbl}
\usepackage{booktabs}
\usepackage{tabularx}
\usepackage{diagbox} % case d'en-tête diagonale des tableaux à double entrée
% Colonnes extensibles centrée / gauche / droite (C, L, R), souvent utilisées par les modèles
\newcolumntype{C}{>{\centering\arraybackslash}X}
\newcolumntype{L}{>{\raggedright\arraybackslash}X}
\newcolumntype{R}{>{\raggedleft\arraybackslash}X}
\usepackage{array}
\usepackage{multicol}
\usepackage{multirow}
\usepackage{siunitx}
% parse-numbers=false : accepte aussi \SI{2,5 \times 10^{-3}}{...}
\sisetup{locale=FR,output-decimal-marker={,},group-minimum-digits=4,parse-numbers=false}
\DeclareSIUnit\ohm{\ensuremath{\symup{\Omega}}} % U+2126 absent de la police
\DeclareSIUnit\division{div} % réglages d'oscilloscope (ms/div, V/div)
\DeclareSIUnit\tr{tr} % tours (vitesse de rotation, ex. \SI{1500}{\tr\per\minute})
\DeclareSIUnit\molar{M} % concentration molaire (ex. \SI{3}{\molar}, \SI{1}{\micro\molar})
\DeclareSIUnit\an{an} % année (le modèle confond parfois avec \year, un compteur TeX) — ex. \SI{5}{\kilo\meter\per\an}
\DeclareSIUnit\FCFA{FCFA} % franc CFA (ex. \SI{500}{\FCFA\per\kilo\gram})
\DeclareSIUnit\degreeBrix{\textdegree Brix} % teneur en sucre d'un sirop/jus (réfractométrie)
\DeclareSIUnit\month{mois} % le modèle utilise parfois \month, un compteur TeX, comme unité de durée
\usepackage{qrcode}
% \degree : commande fréquemment utilisée par les modèles pour un angle ou une
% température en dehors de \SI{}{} (non fournie par les paquets chargés).
\providecommand{\degree}{^{\circ}}
% \qed (fin de démonstration), utilisé par les modèles sans amsthm
\providecommand{\qed}{\hfill$\square$}
% \barre : double barre des valeurs interdites dans un tableau de variations (tkz-tab)
\providecommand{\barre}{\ensuremath{\|}}
% environnement proof (amsthm non chargé) utilisé par les modèles
\ifdefined\proof\else\newenvironment{proof}[1][Démonstration]{\par\noindent\textit{#1.}\ }{\qed\par}\fi
\usepackage{lastpage}
\usepackage{hyperref}
\hypersetup{hidelinks}

% ============================================================
% Palette « Pop » OnBuch+ — mêmes hex que la classe P (plus_ui.dart)
% ============================================================
\definecolor{poporange}{HTML}{F59321}
\definecolor{popdark}{HTML}{E07A0C}
\definecolor{popcream}{HTML}{FFF6E8}
\definecolor{popink}{HTML}{1C1714}
\definecolor{poppurple}{HTML}{7A5AE0}
\definecolor{popgreen}{HTML}{2E9E6A}
\definecolor{poppink}{HTML}{E0457B}
\definecolor{popblue}{HTML}{2F7DE1}
\definecolor{popgold}{HTML}{C98A12}
\definecolor{popmuted}{HTML}{8A7F76}
\definecolor{popline}{HTML}{EADFD2}
\definecolor{popgray}{HTML}{8A7F76}
\colorlet{popgrey}{popgray}
% Alias tolérants (le modèle nomme parfois les couleurs différemment)
\colorlet{popwhite}{white}
\colorlet{popblack}{popink}
\colorlet{popred}{poppink}
\colorlet{popyellow}{popgold}
% Teintes claires (fonds de tableaux/encadrés) — le modèle nomme parfois
% les couleurs avec un suffixe L (ex. popblueL) sans les avoir définies.
\colorlet{poporangeL}{poporange!20}
\colorlet{popdarkL}{popdark!20}
\colorlet{poppurpleL}{poppurple!20}
\colorlet{popgreenL}{popgreen!20}
\colorlet{poppinkL}{poppink!20}
\colorlet{popblueL}{popblue!20}
\colorlet{popgoldL}{popgold!20}
\colorlet{popredL}{popred!20}
\colorlet{popgrayL}{popgray!20}
% Teintes claires pour fonds de tableaux / zones
\colorlet{popblueL}{popblue!10!white}
\colorlet{poporangeL}{poporange!14!white}
\colorlet{popgreenL}{popgreen!12!white}
\colorlet{poppurpleL}{poppurple!12!white}
\colorlet{poppinkL}{poppink!12!white}
\colorlet{popgoldL}{popgold!14!white}

\pagecolor{popcream}

% ============================================================
% Typographie
% ============================================================
\defaultfontfeatures{Ligatures=TeX}
\setmainfont{PlusJakartaSans-Medium.ttf}[Path=../../../fonts/,BoldFont=PlusJakartaSans-Bold.ttf]
% Symboles, API et emojis absents de la police principale : repli sur DejaVu Sans (caractères actifs)
\newfontface\symfont{DejaVuSans.ttf}[Path=../../../fonts/]
\makeatletter
\newcommand\symactive[1]{\catcode#1=13 \begingroup\lccode`\~=#1 \lowercase{\endgroup\def~}{{\symfont\char#1}}}
\newcommand\symblank[1]{\catcode#1=13 \begingroup\lccode`\~=#1 \lowercase{\endgroup\def~}{}}
\newcommand\symhyph[1]{\catcode#1=13 \begingroup\lccode`\~=#1 \lowercase{\endgroup\def~}{\mbox{-}}}
\newcount\symc
\newcommand\symrange[2]{\symc=#1 \@whilenum\symc<#2 \do{\expandafter\symactive\expandafter{\the\symc}\advance\symc by 1 }}
\newcommand\symblankrange[2]{\symc=#1 \@whilenum\symc<#2 \do{\expandafter\symblank\expandafter{\the\symc}\advance\symc by 1 }}
\makeatother
\symrange{"0250}{"02FF}\symactive{"2713}\symactive{"2714}\symactive{"2717}\symactive{"2718}\symactive{"2610}\symactive{"2611}\symactive{"2612}
\symhyph{"2011}\symhyph{"2E1A}\symblankrange{"1F300}{"1FAFF}\symblank{"FE0F}
% Maths en sans-serif (Fira Math) avec chiffres et lettres droites de Plus
% Jakarta, pour que les formules se fondent dans le texte.
\usepackage[math-style=ISO,bold-style=ISO]{unicode-math}
\setmathfont{FiraMath-Regular.otf}[Path=../../../fonts/]
\setmathfont{PlusJakartaSans-Medium.ttf}[Path=../../../fonts/,range={up/num,up/{Latin,latin},\mathup}]
\usepackage{icomma} % virgule décimale sans espace en mode math
% Caractères mathématiques Unicode absents des polices de texte (Plus Jakarta,
% Archivo Black) : rendus par leur équivalent mathématique, en texte comme en
% formule — sinon ils s'affichent en carré vide (ex. « Arithmétique dans ℤ »).
\usepackage{newunicodechar}
\newunicodechar{ℝ}{\ensuremath{\mathbb{R}}}
\newunicodechar{ℤ}{\ensuremath{\mathbb{Z}}}
\newunicodechar{ℂ}{\ensuremath{\mathbb{C}}}
\newunicodechar{ℕ}{\ensuremath{\mathbb{N}}}
\newunicodechar{ℚ}{\ensuremath{\mathbb{Q}}}
\newunicodechar{𝔻}{\ensuremath{\mathbb{D}}}
\newunicodechar{α}{\ensuremath{\alpha}}
\newunicodechar{β}{\ensuremath{\beta}}
\newunicodechar{γ}{\ensuremath{\gamma}}
\newunicodechar{δ}{\ensuremath{\delta}}
\newunicodechar{ε}{\ensuremath{\varepsilon}}
\newunicodechar{θ}{\ensuremath{\theta}}
\newunicodechar{λ}{\ensuremath{\lambda}}
\newunicodechar{μ}{\ensuremath{\mu}}
\newunicodechar{σ}{\ensuremath{\sigma}}
\newunicodechar{τ}{\ensuremath{\tau}}
\newunicodechar{φ}{\ensuremath{\varphi}}
\newunicodechar{ψ}{\ensuremath{\psi}}
\newunicodechar{ω}{\ensuremath{\omega}}
\newunicodechar{Σ}{\ensuremath{\Sigma}}
% majuscules produites par \MakeUppercase dans les titres (α-aminés -> Α)
\newunicodechar{Α}{\ensuremath{\alpha}}
\newunicodechar{Β}{\ensuremath{\beta}}
\newunicodechar{Γ}{\ensuremath{\Gamma}}
\newunicodechar{Δ}{\ensuremath{\Delta}}
\newunicodechar{Ε}{\ensuremath{\varepsilon}}
\newunicodechar{Θ}{\ensuremath{\Theta}}
\newunicodechar{Λ}{\ensuremath{\Lambda}}
\newunicodechar{Μ}{\ensuremath{\mu}}
\newunicodechar{Τ}{\ensuremath{\tau}}
\newunicodechar{Φ}{\ensuremath{\Phi}}
\newunicodechar{Ψ}{\ensuremath{\Psi}}
\newunicodechar{Ω}{\ensuremath{\Omega}}
\newunicodechar{↦}{\ensuremath{\mapsto}}
\newunicodechar{∈}{\ensuremath{\in}}
\newunicodechar{∉}{\ensuremath{\notin}}
\newunicodechar{∧}{\ensuremath{\wedge}}
\newunicodechar{⇔}{\ensuremath{\Leftrightarrow}}
\newunicodechar{⟺}{\ensuremath{\Longleftrightarrow}}
\newunicodechar{⇒}{\ensuremath{\Rightarrow}}
\newunicodechar{⟹}{\ensuremath{\Longrightarrow}}
\newunicodechar{∘}{\ensuremath{\circ}}
\newunicodechar{∪}{\ensuremath{\cup}}
\newunicodechar{∩}{\ensuremath{\cap}}
\newunicodechar{≡}{\ensuremath{\equiv}}
\newunicodechar{⊂}{\ensuremath{\subset}}
\newunicodechar{⊥}{\ensuremath{\perp}}
\newunicodechar{∃}{\ensuremath{\exists}}
\newunicodechar{∀}{\ensuremath{\forall}}
\newunicodechar{∛}{\ensuremath{\sqrt[3]{\,}}}
\newunicodechar{∜}{\ensuremath{\sqrt[4]{\,}}}
\newunicodechar{⃗}{\ensuremath{{}^{\to}}}
\newunicodechar{ⁿ}{\ifmmode^{n}\else\textsuperscript{\ensuremath{n}}\fi}
\newunicodechar{ˣ}{\ifmmode^{x}\else\textsuperscript{\ensuremath{x}}\fi}
\newunicodechar{ᵇ}{\ifmmode^{b}\else\textsuperscript{\ensuremath{b}}\fi}
\newunicodechar{ʳ}{\ifmmode^{r}\else\textsuperscript{\ensuremath{r}}\fi}
\newunicodechar{ᵘ}{\ifmmode^{u}\else\textsuperscript{\ensuremath{u}}\fi}
\newunicodechar{ʸ}{\ifmmode^{y}\else\textsuperscript{\ensuremath{y}}\fi}
\newunicodechar{ᵃ}{\ifmmode^{a}\else\textsuperscript{\ensuremath{a}}\fi}
\newunicodechar{ᵉ}{\ifmmode^{e}\else\textsuperscript{\ensuremath{e}}\fi}
\newunicodechar{ᵏ}{\ifmmode^{k}\else\textsuperscript{\ensuremath{k}}\fi}
\newunicodechar{ᵖ}{\ifmmode^{p}\else\textsuperscript{\ensuremath{p}}\fi}
\newunicodechar{ᴺ}{\ifmmode^{N}\else\textsuperscript{\ensuremath{N}}\fi}
\newunicodechar{ᶜ}{\ifmmode^{c}\else\textsuperscript{\ensuremath{c}}\fi}
\newunicodechar{ʰ}{\ifmmode^{h}\else\textsuperscript{\ensuremath{h}}\fi}
\newunicodechar{ᵠ}{\ifmmode^{q}\else\textsuperscript{\ensuremath{q}}\fi}
\newunicodechar{ₙ}{\ifmmode_{n}\else\textsubscript{\ensuremath{n}}\fi}
\newunicodechar{ₐ}{\ifmmode_{a}\else\textsubscript{\ensuremath{a}}\fi}
\newunicodechar{ᵢ}{\ifmmode_{i}\else\textsubscript{\ensuremath{i}}\fi}
\newunicodechar{ᵣ}{\ifmmode_{r}\else\textsubscript{\ensuremath{r}}\fi}
\newunicodechar{ₖ}{\ifmmode_{k}\else\textsubscript{\ensuremath{k}}\fi}
\newunicodechar{ᵦ}{\ifmmode_{\beta}\else\textsubscript{\ensuremath{\beta}}\fi}
\newunicodechar{ₓ}{\ifmmode_{x}\else\textsubscript{\ensuremath{x}}\fi}
\newfontfamily\popdisplay{ArchivoBlack-Regular.ttf}[Path=../../../fonts/]
\newfontfamily\poplogo{SpaceGrotesk-Bold.ttf}[Path=../../../fonts/]
\newfontfamily\popsemi{PlusJakartaSans-SemiBold.ttf}[Path=../../../fonts/]
\newfontfamily\popxbold{PlusJakartaSans-ExtraBold.ttf}[Path=../../../fonts/]
\newfontfamily\popxboldls{PlusJakartaSans-ExtraBold.ttf}[Path=../../../fonts/,LetterSpace=3.0]

\setlist[enumerate]{leftmargin=1.7em,itemsep=5pt,topsep=4pt}
\setlist[itemize]{leftmargin=1.7em,itemsep=4pt,topsep=4pt,label={\color{poporange}\textbullet}}
\setlength{\parindent}{0pt}
\setlength{\parskip}{5pt}
\renewcommand{\arraystretch}{1.25}

% pgfplots : style Pop par défaut
\pgfplotsset{
  every axis/.append style={axis line style={popink,line width=1pt},tick style={popink},
    grid style={popline,line width=0.5pt},label style={font=\small},tick label style={font=\footnotesize}},
  popcurve/.style={poporange,line width=1.8pt,no markers,smooth},
}
% Même style disponible dans un \draw TikZ simple (hors pgfplots)
\tikzset{popcurve/.style={poporange,line width=1.8pt}}
\tikzset{popgrid/.style={popline,very thin}}
\colorlet{popcurve}{poporange} % le nom du style est parfois utilisé comme couleur

% ============================================================
% Pill / squiggle
% ============================================================
\newcommand{\poppill}[3]{%
  \tikz[baseline=-0.55ex]\node[draw=popink,line width=1.2pt,rounded corners=2.6pt,%
    fill=#2,inner xsep=6.5pt,inner ysep=3pt]{%
    \popxboldls\fontsize{7}{7}\selectfont\color{#3}\MakeUppercase{#1}};%
}
\newcommand{\popsquiggle}[1]{%
  \tikz[baseline=-0.4ex]{
    \draw[#1,line width=2.6pt,line cap=round]
      plot [smooth] coordinates {(0,0.09) (0.34,0.01) (0.68,0.21) (1.02,0.01) (1.36,0.10)};
  }%
}
% Mot-clé surligné (marqueur orange sous le texte)
\newcommand{\cle}[1]{{\popxbold\color{popdark}#1}}

% ============================================================
% En-tête / pied
% ============================================================
\pagestyle{fancy}
\fancyhf{}
\renewcommand{\headrulewidth}{0pt}
\renewcommand{\footrulewidth}{0pt}
\lhead{\raisebox{-0.15ex}{\poplogo\fontsize{13}{13}\selectfont\color{popink}OnBuch\color{poporange}+}}
\rhead{\poppill{\DOCMATIERE\ \textbullet\ \DOCNIVEAU\ \textbullet\ Leçon \DOCLECON}{white}{popink}}
\lfoot{\popxbold\fontsize{7.6}{7.6}\selectfont\color{popmuted}\MakeUppercase{Cours signé OnBuch+}\ \textbullet\ {\popsemi\DOCTITRE}}
\rfoot{\tikz[baseline=-0.6ex]\node[rounded corners=8.5pt,fill=popink,minimum height=17pt,minimum width=17pt,inner xsep=5pt,inner sep=0]{\popxbold\fontsize{8}{8}\selectfont\color{popcream}\thepage\,/\,\pageref*{LastPage}};}

% ============================================================
% Titres
% ============================================================
\newcommand{\chaptitre}[2]{%
  \begin{tcolorbox}[enhanced,colback=poporange,colframe=popink,boxrule=2.6pt,arc=11pt,
      drop shadow=popink,left=15pt,right=15pt,top=12pt,bottom=11pt,before skip=3pt,after skip=18pt]
    \poppill{#2}{popink}{popcream}\par\vspace{9pt}
    {\raggedright\hyphenpenalty=10000\exhyphenpenalty=10000\popdisplay\fontsize{22}{24}\selectfont\color{popcream}\MakeUppercase{#1}\par}
  \end{tcolorbox}
}
% Section numérotée : pastille orange avec le numéro + titre Archivo
\newcounter{popsec}
\newcommand{\coursec}[1]{%
  \stepcounter{popsec}\setcounter{popsub}{0}%
  \par\vspace{16pt}\noindent
  \begin{minipage}{\linewidth}
  \tikz[baseline=-0.7ex]\node[draw=popink,line width=1.6pt,fill=poporange,rounded corners=4pt,minimum size=22pt,inner sep=0]{\popdisplay\fontsize{12}{12}\selectfont\color{popcream}\thepopsec};%
  \hspace{8pt}{\popdisplay\fontsize{15}{17}\selectfont\color{popink}\MakeUppercase{#1}}\par
  \vspace{4pt}\noindent{\color{poporange}\rule{36pt}{3.6pt}}
  \end{minipage}\par\nopagebreak\vspace{8pt}
}
\newcounter{popsub}
\newcommand{\courssub}[1]{%
  \stepcounter{popsub}%
  \par\vspace{9pt}\noindent{\popxbold\fontsize{11.5}{13}\selectfont\color{popink}{\color{poporange}\thepopsec.\thepopsub}\hspace{6pt}#1}\par\nopagebreak\vspace{4pt}
}

% ============================================================
% Boîtes pédagogiques — même recette : fond blanc, bordure épaisse colorée,
% ombre dure encre, badge plein. Titre optionnel [#1].
% ============================================================
\tcbset{popbase/.style={breakable,enhanced,colback=white,boxrule=2.3pt,arc=9pt,
  shadow={3pt}{-3pt}{0pt}{popink},left=14pt,right=14pt,top=11pt,bottom=11pt,before skip=11pt,after skip=15pt,
  fontupper=\popsemi\fontsize{10.6}{14.5}\selectfont\color{popink},
  fontlower=\popsemi\fontsize{10.6}{14.5}\selectfont\color{popink}}}
% Coche dessinée (le glyphe ✓ n'existe pas dans les polices du document)
\newcommand{\popcheck}[1]{\tikz[baseline=-0.5ex]\draw[#1,line width=1.6pt,line cap=round,line join=round](-0.13,0.0)--(-0.04,-0.09)--(0.13,0.1);}
\providecommand{\coche}{\popcheck{popgreen}}
\providecommand{\resultat}[1]{\par\smallskip\noindent\fcolorbox{popgreen}{popgreenL}{\parbox{\dimexpr\linewidth-2\fboxsep-2\fboxrule\relax}{\textbf{Résultat.} #1}}\par\smallskip}
\newcommand{\popboxhead}[3]{\poppill{#1}{#2}{white}\ifx\relax#3\relax\else\hspace{6pt}{\popxbold\fontsize{10.6}{12}\selectfont\color{popink}#3}\fi\par\vspace{7pt}}

\newtcolorbox{aretenir}[1][]{popbase,colframe=popgreen,before upper={\popboxhead{À retenir}{popgreen}{#1}}}
% Texte en espagnol (document de lecture, dialogue, poème, extrait) : cadre
% d'étude. Un titre optionnel (source, auteur) s'affiche dans l'en-tête.
\newtcolorbox{textebox}[1][]{popbase,colframe=popblue,before upper={\popboxhead{Texto}{popblue}{#1}\begin{otherlanguage}{spanish}},after upper={\end{otherlanguage}}}
% Marques ✓ / ✗ (forme correcte / fautive) souvent utilisées par les modèles
\providecommand{\cross}{\textcolor{poppink}{\ensuremath{\times}}}
\providecommand{\xmark}{\cross}\providecommand{\crossmark}{\cross}\providecommand{\cmark}{\ensuremath{\checkmark}}
% Ligne de réponse / justification à compléter (inventée par les modèles)
\providecommand{\justification}[1]{\par\noindent\textit{Justification~:} \dotfill\par}
% \textipa (package tipa non chargé) : la police ne couvre pas l'API -> simple passage
\providecommand{\textipa}[1]{#1}
% Soulignement (ulem non chargé) : \uline, \ul
\providecommand{\uline}[1]{\underline{#1}}\providecommand{\ul}[1]{\underline{#1}}
% \glossaire{\item ...} inventée par les modèles : liste de glossaire
\providecommand{\glossaire}[1]{\begin{itemize}#1\end{itemize}}
% \alert (beamer) inventée par les modèles : texte mis en évidence
\providecommand{\alert}[1]{\textbf{\textcolor{poppink}{#1}}}
% variantes ASCII d'ordinaux écrites en macro
\providecommand{\iere}{\textsuperscript{re}}\providecommand{\ieme}{\textsuperscript{e}}\providecommand{\ere}{\textsuperscript{re}}\providecommand{\eme}{\textsuperscript{e}}\providecommand{\er}{\textsuperscript{er}}\providecommand{\ie}{\textsuperscript{e}}
% Ordinaux français écrits en macro par les modèles : 1\ière, 3\ième
\newcommand{\ière}{\textsuperscript{re}}\newcommand{\ième}{\textsuperscript{e}}
% \exemple (inventée par les modèles) : étiquette « Exemple : »
\providecommand{\exemple}{\textbf{Exemple~:}\ }
% Mot ou phrase en espagnol à l'intérieur d'un paragraphe français
\newcommand{\esp}[1]{\ifmmode\text{\textit{#1}}\else\begingroup\selectlanguage{spanish}\itshape #1\endgroup\fi}
\newtcolorbox{exemplebox}[1][]{popbase,colframe=poporange,before upper={\popboxhead{Exemple}{poporange}{#1}}}
\newtcolorbox{definition}[1][]{popbase,colframe=popblue,before upper={\popboxhead{Définition}{popblue}{#1}}}
\newtcolorbox{propriete}[1][]{popbase,colframe=poppurple,before upper={\popboxhead{Propriété}{poppurple}{#1}}}
\newtcolorbox{methode}[1][]{popbase,colframe=popgold,before upper={\popboxhead{Méthode}{popgold}{#1}}}
\newtcolorbox{attention}[1][]{popbase,colframe=poppink,before upper={\popboxhead{Attention}{poppink}{#1}}}
\newtcolorbox{experience}[1][]{popbase,colframe=popblue,colback=popblueL,before upper={\popboxhead{Expérience}{popblue}{#1}}}
% « Le savais-tu ? » : carte violette pleine (contexte, culture, Cameroun)
\newtcolorbox{savaistu}[1][]{popbase,colback=poppurple,colframe=popink,
  fontupper=\popsemi\fontsize{10.4}{14.2}\selectfont\color{white},
  before upper={\poppill{Le savais-tu\,?}{white}{poppurple}\ifx\relax#1\relax\else\hspace{6pt}{\popxbold\color{white}#1}\fi\par\vspace{7pt}}}
% Exercice résolu : énoncé en haut, \tcblower puis la solution sur fond crème
\newcounter{popexo}\newcounter{popexr}
\newtcolorbox{exoresolu}[1][]{popbase,colframe=poporange,colbacklower=popcream,
  segmentation style={popink,line width=1.2pt,dashed},
  before upper={\stepcounter{popexr}\popboxhead{Exercice résolu \thepopexr}{poporange}{#1}},
  before lower={\poppill{Solution}{popink}{popcream}\par\vspace{6pt}}}
% Exercice (non résolu en place) — la correction est dans la partie Corrigés
\newtcolorbox{exercice}[1][]{popbase,colframe=popink,
  before upper={\stepcounter{popexo}\popboxhead{Exercice \thepopexo}{popink}{#1}}}
% Corrigé
\newtcolorbox{corrige}[1][]{popbase,colframe=popgreen,colback=popgreenL,
  before upper={\popboxhead{Corrigé}{popgreen}{#1}}}
% Objectifs / prérequis / bilan
\newtcolorbox{objectifs}{popbase,colframe=popink,colback=poporangeL,
  before upper={\popboxhead{Objectifs de la leçon}{popink}{}}}
\newtcolorbox{prerequis}{popbase,colframe=popink,colback=popblueL,
  before upper={\popboxhead{Prérequis}{popblue}{}}}
\newtcolorbox{bilan}{popbase,colframe=popink,colback=white,boxrule=2.8pt,
  before upper={\popboxhead{Fiche bilan}{poporange}{L'essentiel à connaître pour l'examen}}}
% Figure encadrée avec légende
\newtcolorbox{popfigure}[1][]{enhanced,breakable=false,colback=white,colframe=popline,boxrule=1.4pt,arc=8pt,
  left=10pt,right=10pt,top=10pt,bottom=8pt,before skip=10pt,after skip=12pt,halign=center,
  lower separated=false,halign lower=center,
  fontlower=\popsemi\fontsize{9}{11.5}\selectfont\color{popmuted},
  #1}
\newcommand{\legende}[1]{\par\vspace{6pt}{\popsemi\fontsize{9}{11.5}\selectfont\color{popmuted}#1}}

% ============================================================
% Signature OnBuch+
% ============================================================
\newcommand{\poplogoinline}[1]{{\poplogo\fontsize{#1}{#1}\selectfont\color{popink}OnBuch\color{poporange}+}}
\newcommand{\signatureonbuch}{%
  \begin{tcolorbox}[enhanced,colback=popink,colframe=popink,boxrule=2.2pt,arc=10pt,
      drop shadow=poporange,left=14pt,right=14pt,top=10pt,bottom=10pt,before skip=0pt,after skip=0pt]
    \tikz[baseline=-0.7ex]\node[circle,fill=popgreen,draw=popcream,line width=1.3pt,minimum size=22pt,inner sep=0]{\popcheck{white}};%
    \hspace{9pt}\begin{minipage}[c]{0.62\linewidth}
      {\popxbold\fontsize{12}{13}\selectfont\color{popcream}Signé\ \textbullet\ {\poplogo OnBuch\color{poporange}+}}\par\vspace{2pt}
      {\raggedright\popsemi\fontsize{8}{10.5}\selectfont\color{popline}\MakeUppercase{Équipe pédagogique OnBuch+}\\\MakeUppercase{Conforme au programme officiel MINESEC}\par}
    \end{minipage}\hfill
    {\poplogo\fontsize{22}{22}\selectfont\color{popcream}OnBuch\color{poporange}+}
  \end{tcolorbox}%
}

% ============================================================
% Page de garde
% #1 titre · #2 matière · #3 niveau · #4 module · #5 n° leçon · #6 durée ·
% #7 description / objectifs (texte ou itemize)
% ============================================================
\newcommand{\pagegarde}[7]{%
  \thispagestyle{empty}
  \newgeometry{a4paper,margin=1.7cm,top=1.6cm,bottom=1.4cm}%
  \begin{tikzpicture}[remember picture,overlay]
    \fill[poporange!18] ([xshift=-3.2cm,yshift=-3.6cm]current page.north east) circle (4.6cm);
    \fill[poppurple!14] ([xshift=2.4cm,yshift=7.6cm]current page.south west) circle (3.8cm);
  \end{tikzpicture}%
  {\poplogo\fontsize{30}{30}\selectfont\color{popink}OnBuch\color{poporange}+}\hfill
  \raisebox{0.6ex}{\poppill{Cours signé \textbullet\ Premium}{popink}{popcream}}\par
  \vspace{1pt}\popsquiggle{poppurple}\par
  \vspace{8pt}
  \begin{tcolorbox}[enhanced,colback=poporange,colframe=popink,boxrule=2.8pt,arc=13pt,
      drop shadow=popink,left=18pt,right=18pt,top=12pt,bottom=13pt,after skip=10pt]
    \poppill{#2 \textbullet\ #3}{popink}{popcream}\hspace{5pt}\poppill{#4}{white}{popink}\par\vspace{8pt}
    {\popdisplay\fontsize{14}{14}\selectfont\color{popink}LEÇON #5}\par\vspace{4pt}
    {\raggedright\hyphenpenalty=10000\exhyphenpenalty=10000\popdisplay\fontsize{25}{27}\selectfont\color{popcream}\MakeUppercase{#1}\par}
    \vspace{8pt}
    {\popxbold\fontsize{9.5}{9.5}\selectfont\color{popink}\MakeUppercase{Durée conseillée : #6}}
  \end{tcolorbox}
  \begin{tcolorbox}[enhanced,colback=white,colframe=popink,boxrule=2.2pt,arc=10pt,
      drop shadow=popink,left=16pt,right=16pt,top=10pt,bottom=10pt,after skip=8pt]
    \poppill{Dans cette leçon}{poppurple}{white}\par\vspace{5pt}
    \popsemi\fontsize{9.3}{12.6}\selectfont\color{popink}#7
  \end{tcolorbox}
  \noindent\begin{minipage}[t]{0.487\linewidth}
    \begin{tcolorbox}[enhanced,colback=popgreen,colframe=popink,boxrule=2.2pt,arc=10pt,
        drop shadow=popink,left=11pt,right=11pt,top=10pt,bottom=10pt,height=3.1cm,valign=center]
      \tcbox[colback=white,colframe=popink,boxrule=1.3pt,arc=5pt,boxsep=0pt,nobeforeafter,
        left=3pt,right=3pt,top=3pt,bottom=3pt,valign=center]{\qrcode[height=1.9cm]{https://play.google.com/store/apps/details?id=com.luvvix.onbuch&hl=fr}}%
      \hspace{8pt}\begin{minipage}[c]{0.5\linewidth}
        {\popdisplay\fontsize{11}{12.5}\selectfont\color{white}\MakeUppercase{Télécharge OnBuch}}\par\vspace{4pt}
        {\raggedright\popsemi\fontsize{8.4}{11}\selectfont\color{white}Gratuit \textbullet\ Google Play\\Tuteur IA, annales, concours, résultats\par}
      \end{minipage}
    \end{tcolorbox}
  \end{minipage}\hfill
  \begin{minipage}[t]{0.487\linewidth}
    \begin{tcolorbox}[enhanced,colback=poppurple,colframe=popink,boxrule=2.2pt,arc=10pt,
        drop shadow=popink,left=11pt,right=11pt,top=10pt,bottom=10pt,height=3.1cm,valign=center]
      \tikz[baseline=-0.7ex]\node[circle,fill=white,draw=popink,line width=1.3pt,minimum size=1.6cm,inner sep=0]{\popdisplay\fontsize{24}{24}\selectfont\color{poppurple}+};%
      \hspace{8pt}\begin{minipage}[c]{0.6\linewidth}
        {\popdisplay\fontsize{11}{12.5}\selectfont\color{white}\MakeUppercase{Passe à OnBuch+}}\par\vspace{4pt}
        {\raggedright\popsemi\fontsize{8.4}{11}\selectfont\color{white}Tous les cours signés, Tuteur IA illimité, fiches PDF — onglet OnBuch+ de l'app\par}
      \end{minipage}
    \end{tcolorbox}
  \end{minipage}
  \vspace{8pt}
  \popcontenu
  \vfill
  \signatureonbuch
  \restoregeometry
  \newpage
}

% Rangée « Ce cours contient » (page de garde)
\newcommand{\poptile}[3]{% #1 couleur #2 gros texte #3 légende
  \begin{tcolorbox}[enhanced,colback=#1,colframe=popink,boxrule=2pt,arc=9pt,shadow={2.5pt}{-2.5pt}{0pt}{popink},
      left=6pt,right=6pt,top=9pt,bottom=9pt,halign=center,valign=center,height=1.75cm,width=\linewidth,nobeforeafter]
    {\popdisplay\fontsize{12.5}{13}\selectfont\color{white}\MakeUppercase{#2}}\par\vspace{3pt}
    {\centering\popsemi\fontsize{7.8}{9.5}\selectfont\color{white}#3\par}
  \end{tcolorbox}}
\newcommand{\popcontenu}{%
  \par\noindent{\popxboldls\fontsize{7.5}{7.5}\selectfont\color{popmuted}CE COURS SIGNÉ CONTIENT}\par\vspace{7pt}
  \noindent\begin{minipage}[t]{0.235\linewidth}\poptile{popblue}{Cours}{complet, illustré, pas à pas}\end{minipage}\hfill
  \begin{minipage}[t]{0.235\linewidth}\poptile{poporange}{Méthodes}{et exercices résolus}\end{minipage}\hfill
  \begin{minipage}[t]{0.235\linewidth}\poptile{poppink}{Exercices}{type examen + corrigés}\end{minipage}\hfill
  \begin{minipage}[t]{0.235\linewidth}\poptile{popgreen}{Fiche bilan}{l'essentiel à retenir}\end{minipage}\par}

% Sommaire
\newtcolorbox{sommaire}{enhanced,colback=white,colframe=popink,boxrule=2.2pt,arc=10pt,
  drop shadow=popink,left=18pt,right=18pt,top=13pt,bottom=13pt,before skip=6pt,after skip=20pt,
  before upper={\poppill{Sommaire}{poppurple}{white}\par\vspace{9pt}\popsemi\fontsize{10.6}{14.5}\selectfont\color{popink}}}

% Signature de fin de cours — poussée tout en bas de la dernière page
\newcommand{\findecours}{%
  \par\vfill\nopagebreak
  \begin{center}\popsquiggle{poporange}\end{center}\vspace{4pt}
  \signatureonbuch
}
\AtBeginDocument{\def\llbracket{\mathopen{[\mkern-3.5mu[}}\def\rrbracket{\mathclose{]\mkern-3.5mu]}}} % intervalles d'entiers (glyphe ⟦ absent de la police)
% Glyphes absents de Fira Math : redéfinis à partir de glyphes disponibles
\AtBeginDocument{%
  \def\setminus{\mathbin{\backslash}}\let\smallsetminus\setminus
  \def\vdots{\mathord{\vbox{\baselineskip3.2pt\lineskiplimit0pt\kern4pt\hbox{.}\hbox{.}\hbox{.}}}}%
  \def\ddots{\mathinner{\mkern1mu\raise6.5pt\hbox{.}\mkern2mu\raise3.6pt\hbox{.}\mkern2mu\raise0.7pt\hbox{.}\mkern1mu}}%
  \def\triangle{\mathord{\scalebox{0.9}{$\symup{\Delta}$}}}%
  \def\nabla{\mathord{\raisebox{\depth}{\scalebox{1}[-1]{$\symup{\Delta}$}}}}%
  \def\checkmark{\popcheck{popdark}}%
  \def\star{\mathbin{\ast}}\def\bigstar{\mathord{\ast}}%
  \def\diamond{\mathbin{\scalebox{0.6}{$\Diamond$}}}%
  \def\bigcup{\mathop{\mathchoice{\raisebox{-0.3em}{\scalebox{1.7}{$\cup$}}}{\scalebox{1.25}{$\cup$}}{\cup}{\cup}}\displaylimits}%
  \def\bigcap{\mathop{\mathchoice{\raisebox{-0.3em}{\scalebox{1.7}{$\cap$}}}{\scalebox{1.25}{$\cap$}}{\cap}{\cap}}\displaylimits}%
  \def\lesseqqgtr{\mathrel{\lessgtr}}\def\bigoplus{\mathop{\oplus}\displaylimits}%
}
\providecommand{\ssi}{si et seulement si} % abréviation fréquente
\AtBeginDocument{\providecommand{\wideparen}{\overparen}} % arc de cercle

%% FIGFIX-BEGIN v5 — correctifs globaux des figures (docs/onbuch-generation/tools/figfix)
\usepackage{adjustbox}\usepackage{etoolbox}\tcbuselibrary{raster}
% toute figure TikZ/pgfplots plus large que la ligne est réduite à la largeur disponible
\BeforeBeginEnvironment{tikzpicture}{\begin{adjustbox}{max width=\linewidth}}
\AfterEndEnvironment{tikzpicture}{\end{adjustbox}}
% légendes pgfplots lisibles par-dessus les courbes
\pgfplotsset{every axis/.append style={legend style={fill=white,fill opacity=0.92,text opacity=1,draw=black!15,inner sep=3pt}}}
% étiquettes d'arêtes (syntaxe edge["..."]) avec fond blanc
\tikzset{every edge quotes/.append style={fill=white,inner sep=1.2pt}}
%% FIGFIX-END
\begin{document}

\pagegarde{Traduction : « avoir beau » ; « demander » ; « depuis »}{Espagnol}{Tle A}{Módulo 5 : Médias, communication et culture de la paix}{E25}{2 h}{Cette leçon de traduction (2 h) permet de maîtriser les équivalents espagnols de « avoir beau », « demander » et « depuis » à travers l'observation, la règle, des tableaux comparatifs, des exercices de thème et de version, et une activité d'intégration ancrée dans le contexte camerounais et hispanophone.
\vspace{4pt}
\begin{multicols}{2}\small
\begin{itemize}
\item À la fin de cette leçon, je sais traduire « avoir beau » par por más que / por mucho que + subjonctif (concession réelle) et par aunque + indicatif (concession factuelle).
\item Je distingue pedir (demander un objet/service), preguntar (demander une information), rogar (supplier) et solicitar (demander formellement).
\item Je choisis correctement entre desde (point de départ), desde hace / hace ... que (durée avec action encore en cours) et desde que (antériorité).
\item J'identifie le piège du subjonctif après por más que vs l'indicatif après aunque quand le fait est certain.
\item Je construis des phrases de thème (français → espagnol) et de version (espagnol → français) mobilisant les trois notions.
\item Je repère les faux-amis : demander ≠ demandar, depuis ≠ desde toujours.
\end{itemize}\end{multicols}}

\begin{sommaire}
\begin{multicols}{2}
\begin{enumerate}
\item 1. Observation guidée des trois notions
\item 2. Traduire « avoir beau » : concession et insistance
\item 3. Traduire « demander » : cuatro verbos, cuatro matices
\item 4. Traduire « depuis » : repères temporels et aspect
\item 5. Synthèse comparative, pièges et entraînement thème/version
\item Activité d'intégration
\item Méthodes et exercices résolus
\item Exercices
\item Corrigés des exercices
\item Fiche bilan
\end{enumerate}
\end{multicols}
\end{sommaire}

\begin{objectifs}
\begin{itemize}
\item À la fin de cette leçon, je sais traduire « avoir beau » par por más que / por mucho que + subjonctif (concession réelle) et par aunque + indicatif (concession factuelle).
\item Je distingue pedir (demander un objet/service), preguntar (demander une information), rogar (supplier) et solicitar (demander formellement).
\item Je choisis correctement entre desde (point de départ), desde hace / hace ... que (durée avec action encore en cours) et desde que (antériorité).
\item J'identifie le piège du subjonctif après por más que vs l'indicatif après aunque quand le fait est certain.
\item Je construis des phrases de thème (français → espagnol) et de version (espagnol → français) mobilisant les trois notions.
\item Je repère les faux-amis : demander ≠ demandar, depuis ≠ desde toujours.
\item J'applique la règle de concordance des temps avec desde que (indicatif ou subjonctif selon l'antériorité).
\item Je produis un court texte cohérent (dialogue, mail) intégrant les trois structures.
\end{itemize}
\end{objectifs}

\begin{prerequis}
\begin{itemize}
\item Formation du subjonctif présent (réguliers et irréguliers fréquents).
\item Emploi de l'indicatif présent, passé composé et imparfait pour exprimer l'antériorité et la durée.
\item Différence de base entre pedir et preguntar (vue en Première).
\item Notion de concession avec aunque + indicatif/subjonctif (Première).
\item Vocabulaire de base sur les médias, la communication et la culture de la paix (module 5).
\end{itemize}
\end{prerequis}

\newpage

\coursec{Observation guidée des trois notions}

Au marché central de Yaoundé, un touriste espagnol demande : « ¿Desde cuándo vende usted mangues ? » Le vendeur répond : « Desde hace diez años. » Plus tard, le touriste insiste : « Tengo que preguntar el precio, pero por más que pregunto, no me contesta. » Ces trois situations — durée, question, concession — illustrent les pièges de traduction que nous allons lever aujourd'hui.

\courssub{Corpus d'observation : six situations authentiques}

\begin{textebox}[Corpus d'observation — Médias et vie quotidienne]
\textbf{1.} (Presse, \emph{El País}, société) \esp{Por más que el gobierno anuncia reformas, la desigualdad persiste en las zonas rurales.}
\textbf{2.} (Dialogue, marché de Douala) \esp{— ¿Desde cuándo trabaja usted aquí? — Desde hace quince años, señor.}
\textbf{3.} (Reportage, \emph{BBC Mundo}, technologie) \esp{Los usuarios pidieron una aplicación más segura y la empresa lanzó una actualización.}
\textbf{4.} (Entretien, radio \emph{RTVGE}, Malabo) \esp{El ministro preguntó por qué no se habían entregado los libros a tiempo.}
\textbf{5.} (Littérature, roman contemporain) \esp{Aunque estudió mucho, no aprobó el examen de selectividad.}
\textbf{6.} (Administratif, communiqué officiel) \esp{La asociación solicitó formalmente una reunión con el alcalde.}
\end{textebox}

\noindent\textbf{Glossaire marginal :}
\begin{itemize}
    \item \esp{anunciar} : annoncer ; \esp{desigualdad} : inégalité ; \esp{persistir} : persister.
    \item \esp{lanzar} : lancer ; \esp{actualización} : mise à jour.
    \item \esp{entregar} : remettre, livrer ; \esp{a tiempo} : à temps.
    \item \esp{aprobar} : réussir (un examen) ; \esp{selectividad} : examen d'entrée à l'université (Espagne).
    \item \esp{formalmente} : formellement ; \esp{alcalde} : maire.
\end{itemize}

\courssub{Repérage visuel et classification intuitive}

Observez les segments en \textbf{gras} dans le corpus. Ils correspondent aux trois notions cibles de cette leçon.

\begin{popfigure}
\begin{tikzpicture}[
    node distance=1.2cm and 1.5cm,
    box/.style={draw=popblue, thick, rounded corners, fill=popblueL, align=center, minimum width=3.5cm, minimum height=1.8cm},
    arrow/.style={->, thick, popdark}
]
\node[box, align=center] (obs){Observation \\ 6 phrases authentiques \\ → Repérage des segments};
\node[box, below left=of obs, align=center] (cat1){Concession \\ \textbf{Por más que} + subj. (phr. 1) \\ \textbf{Aunque} + ind. (phr. 5)};
\node[box, below=of obs, align=center] (cat2){Demande \\ \textbf{Pidieron} (objet, phr. 3) \\ \textbf{Preguntó} (info, phr. 4) \\ \textbf{Solicitó} (formel, phr. 6)};
\node[box, below right=of obs, align=center] (cat3){Temporalité \\ \textbf{Desde cuándo / Desde hace} (phr. 2)};
\draw[arrow] (obs) -- (cat1);
\draw[arrow] (obs) -- (cat2);
\draw[arrow] (obs) -- (cat3);
\end{tikzpicture}
\legende{Classification intuitive des six énoncés du corpus}
\end{popfigure}

\begin{definition}[Concession]
Relation logique où un fait (la cause) n'empêche pas la réalisation d'un autre fait (la conséquence). En français : « avoir beau + infinitif », « quelque ... que », « bien que ». En espagnol : \esp{por más que / por mucho que / aunque} + \textbf{subjonctif} (si hypothétique/non réalisé) ou \textbf{indicatif} (fait réel, admis).
\end{definition}

\begin{definition}[Demande]
Acte de langage visant à obtenir quelque chose. Trois nuances majeures :
\begin{enumerate}
    \item \textbf{Objet / action} : \esp{pedir} (demander un objet, un service, une action).
    \item \textbf{Information} : \esp{preguntar} (poser une question).
    \item \textbf{Formalité / instance supérieure} : \esp{solicitar} (demande écrite, administrative) ; \esp{rogar} (supplier, insister avec politesse).
\end{enumerate}
\end{definition}

\begin{definition}[Depuis]
Repère temporel marquant le point de départ (\emph{terminus a quo}) d'un processus.
\begin{itemize}
    \item \textbf{Action encore en cours} (présent) : \esp{desde hace} + durée / \esp{hace} + durée + \esp{que} + présent / \esp{desde} + date/heure.
    \item \textbf{Action passée terminée} (passé) : \esp{desde hace} + durée / \esp{hace} + durée + \esp{que} + passé / \esp{desde} + date.
    \item \textbf{Antériorité pure (subordonnée)} : \esp{desde que} + indicatif (fait réel) / subjonctif (futur, hypothétique).
\end{itemize}
\end{definition}

\courssub{Formulation d'hypothèses de traduction}

Avant l'apport théorique complet (sections 2 à 4), formulons des hypothèses à partir du corpus.

\begin{center}
\begin{tabularx}{\linewidth}{|X|X|X|X|}
\hline
\rowcolor{popblueL}
\textbf{Phrase (n°)} & \textbf{Segment français cible} & \textbf{Hypothèse espagnol (forme observée)} & \textbf{Contexte / Registre} \\
\hline
1 & \emph{Avoir beau} annoncer & \esp{Por más que} + \textbf{subjonctif} (\esp{anuncia}) & Presse, constat d'échec \\
\hline
2 & \emph{Depuis} quand / \emph{depuis} 15 ans & \esp{Desde cuándo} / \esp{Desde hace} + durée & Dialogue oral, durée en cours \\
\hline
3 & \emph{Demander} (une appli) & \esp{Pedir} (prétérit) & Journalisme, action concrète \\
\hline
4 & \emph{Demander} (pourquoi) & \esp{Preguntar} (prétérit) & Oral / reportage, recherche info \\
\hline
5 & \emph{Bien qu'}il ait étudié / \emph{Avoir beau} étudier & \esp{Aunque} + \textbf{indicatif} (\esp{estudió}) & Récit, fait réel admis \\
\hline
6 & \emph{Demander} (formellement) & \esp{Solicitar} (prétérit) & Administratif, écrit officiel \\
\hline
\end{tabularx}
\end{center}

\courssub{Variations selon le contexte : registre et certitude}

\begin{propriete}[Variation concession : indicatif vs subjonctif]
\esp{Aunque} admet les deux modes selon la valeur de vérité accordée au fait :
\begin{itemize}
    \item \textbf{Indicatif} : le fait est \textbf{réel, certain, admis} par le locuteur. \esp{Aunque estudió mucho, no aprobó.} (Il a vraiment beaucoup étudié).
    \item \textbf{Subjonctif} : le fait est \textbf{hypothétique, non réalisé, mis en doute} ou présenté comme une concession générale. \esp{Aunque estudie mucho, no aprobará.} (Qu'il étudie beaucoup ou pas, il ne réussira pas).
\end{itemize}
\esp{Por más que / por mucho que} exigent presque toujours le \textbf{subjonctif} (valeur de « quelque ... que »).
\end{propriete}

\begin{propriete}[Variation demande : transitivité et registre]
\begin{center}
\begin{tabular}{|l|l|l|l|}
\hline
\rowcolor{popblueL}
\textbf{Verbe} & \textbf{Construction} & \textbf{Objet} & \textbf{Registre} \\
\hline
\esp{Pedir} & \esp{pedir algo (a algn)} & Chose, service, action & Courant \\
\hline
\esp{Preguntar} & \esp{preguntar algo / por algo} & Information & Courant \\
\hline
\esp{Solicitar} & \esp{solicitar algo (a algn)} & Droit, document, audience & Formel, écrit \\
\hline
\esp{Rogar} & \esp{rogar a algn que + subj.} & Action d'autrui & Solennel, poli \\
\hline
\end{tabular}
\end{center}
\end{propriete}

\begin{methode}[Observer pour traduire : 3 étapes]
\begin{enumerate}
    \item \textbf{Identifier la notion en français} : Concession (\emph{avoir beau, bien que}) ? Demande (\emph{demander, demander à, prier}) ? Temporalité (\emph{depuis, il y a ... que}) ?
    \item \textbf{Choisir l'outil espagnol selon le sens précis} :
    \begin{itemize}
        \item Concession réelle $\rightarrow$ \esp{que} + indicatif ; concession hypothétique $\rightarrow$ \esp{que} + subjonctif / \esp{por más que} + subj.
        \item Demande objet $\rightarrow$ \esp{pedir} ; info $\rightarrow$ \esp{preguntar} ; formel $\rightarrow$ \esp{solicitar} ; insistant/poli $\rightarrow$ \esp{rogar}.
        \item Durée en cours $\rightarrow$ \esp{desde hace} / \esp{hace ... que} + présent ; point de départ daté $\rightarrow$ \esp{desde} + date ; subordonnée $\rightarrow$ \esp{desde que}.
    \end{itemize}
    \item \textbf{Vérifier la structure} : Mode (ind./subj.), préposition (\esp{a}, \esp{por}), temps (présent/prétérit/imparfait selon aspect).
\end{enumerate}
\end{methode}

\begin{attention}[Faux ami majeur : \texorpdfstring{\esp{demandar}}{demandar} $\neq$ \emph{demander}]
\esp{Demandar} signifie « \textbf{exiger, réclamer (juridiquement), porter plainte} ».
\begin{itemize}
    \item \esp{El sindicato demandó a la empresa por despido injustificado.} (Le syndicat a \textbf{assigné / porté plainte contre} l'entreprise pour licenciement abusif).
    \item Pour traduire « demander » (standard) : \esp{pedir, preguntar, solicitar, rogar}. \textbf{Jamais \esp{demandar}} sauf contexte juridique.
\end{itemize}
\end{attention}

\begin{exoresolu}[Analyse contrastée : \emph{« Il a beau dire, je ne le crois pas. »}]
\textbf{Énoncé :} Traduire la phrase en espagnol en justifiant le choix du mode.
\tcblower
\textbf{Solution détaillée :}
\begin{enumerate}
    \item \textbf{Notion} : Concession (\emph{avoir beau} + infinitif). Le fait (« il dit ») est réel (il parle), mais le locuteur le présente comme \textbf{inefficace} pour changer son opinion. C'est une concession \emph{adversative} forte.
    \item \textbf{Choix} : \esp{Por más que} (ou \esp{por mucho que}) + \textbf{subjonctif} (valeur de « quelque chose qu'il dise »). \esp{Aunque} + indicatif possible si on admet le fait (« bien qu'il dise »), mais \esp{por más que} rend mieux l'insistance sur la répétition/l'inutilité (\emph{avoir beau}).
    \item \textbf{Traduction} : \esp{Por más que diga, no le creo.} (\esp{diga} = présent subjonctif de \esp{decir}).
    \item \textbf{Variante (fait admis)} : \esp{Aunque dice la verdad, no le creo.} (Indicatif \esp{dice} : il dit la vérité, c'est un fait, mais je ne le crois quand même pas).
\end{enumerate}
\end{exoresolu}

\begin{savaistu}[Le \texorpdfstring{\esp{desde}}{desde} camerounais et la Guinée équatoriale]
En Guinée équatoriale (seul pays africain hispanophone officiel), l'usage de \esp{desde hace} + présent pour exprimer une durée en cours est la norme standard, tout comme en Espagne ou au Mexique. Cependant, à l'oral \emph{popular} (et parfois au Cameroun chez les apprenants), on entend l'anglicisme calqué \esp{desde ... atrás} (\emph{depuis ... en arrière}) ou la confusion \esp{desde dos años} (sans \esp{hace}). \textbf{Règle d'or} : \cle{Durée en cours = \esp{desde hace} + présent} (ou \esp{hace ... que} + présent). Ex. \esp{Vivo en Yaoundé desde hace tres años.} (J'habite Yaoundé depuis trois ans).
\end{savaistu}

\coursec{Traduire « avoir beau » : concession et insistance}

Dans la section précédente, l'observation guidée nous a permis de repérer, dans des phrases authentiques, trois tournures françaises qui n'ont pas d'équivalent mot à mot en espagnol : \emph{avoir beau}, \emph{demander} et \emph{depuis}. Nous abordons maintenant la première en profondeur. « Avoir beau » est l'une des tournures les plus typiquement françaises : l'espagnol ne possède aucun verbe équivalent, et il faut donc raisonner en termes de \cle{sens} — concession, effort vain, fait avéré — avant de choisir la structure espagnole adéquate.

\courssub{Observation : que dit vraiment « avoir beau » ?}

Lisons d'abord ce court texte de presse, rédigé dans un contexte qui nous est familier.

\begin{textebox}[Crónica deportiva — texto original]
El equipo de Douala entrena cada día bajo el sol. Por más que corran los delanteros, el balón no entra en la portería. El entrenador, por mucho que grite desde la banda, no consigue cambiar el resultado. Los aficionados tienen paciencia, pero por fuerte que sea su apoyo, la victoria se resiste. Aunque el capitán juega con el corazón, los errores se repiten partido tras partido. Algunos periodistas escriben que el equipo tiene fácil clasificarse si gana los dos próximos encuentros, pero los jugadores saben que nada será sencillo. « Por más que lo intentamos, la suerte nos abandona », declaró el portero al final del partido. Sin embargo, el presidente del club mantiene la calma: aunque lluevan críticas, el proyecto deportivo continuará. La afición, fiel como siempre, volverá al estadio el domingo que viene.
\end{textebox}

\textbf{Glossaire :} \esp{la banda} : le bord du terrain ; \esp{se resiste} : se fait attendre, résiste ; \esp{clasificarse} : se qualifier ; \esp{el portero} : le gardien de but ; \esp{llover críticas} : pleuvoir des critiques ; \esp{la afición} : les supporters.

\textbf{Questions de compréhension :}
\begin{enumerate}
\item Relevez les quatre constructions différentes qui expriment la concession.
\item Dans \esp{Por más que corran los delanteros}, le verbe est-il au subjonctif ou à l'indicatif ? Pourquoi ce choix de mode ?
\item Quelle tournure du texte traduit « avoir beau jeu de » ?
\end{enumerate}

\courssub{La structure de base : \esp{por más que} / \esp{por mucho que} + subjonctif}

\begin{propriete}[Règle fondamentale]
\esp{Por más que} et \esp{por mucho que} exigent \textbf{toujours le subjonctif} dans l'espagnol standard, car ils expriment une concession \emph{non réalisée} : l'effort est réel, mais le résultat espéré ne se produit pas. C'est l'équivalent direct de « avoir beau » + infinitif.
\begin{center}
\esp{Por más que} / \esp{Por mucho que} + \textbf{subjonctif} , (résultat négatif ou contraire)
\end{center}
Au passé, on emploie l'imparfait du subjonctif : \esp{por más que estudiara/estudiase}.
\end{propriete}

\begin{exemplebox}[Exemples traduits]
\esp{Por más que estudie, no aprueba.} = Il a beau étudier, il ne réussit pas.\par
\esp{Por mucho que grites, nadie te oye.} = Tu as beau crier, personne ne t'entend.\par
\esp{Por más que corriera, llegó tarde al examen.} = Il eut beau courir, il arriva en retard à l'examen.\par
\esp{Por mucho que insistamos, el profesor no cambia de opinión.} = Nous avons beau insister, le professeur ne change pas d'avis.
\end{exemplebox}

Remarquez la nuance : \esp{por mucho que} insiste sur la \cle{quantité} de l'effort (« beaucoup »), \esp{por más que} sur son \cle{intensité} (« davantage »). Dans la pratique, les deux sont interchangeables pour traduire « avoir beau ».

\begin{exemplebox}[Une autre tournure utile : \esp{de nada sirve que}]
Pour insister sur l'\textbf{inutilité} de l'effort, l'espagnol offre aussi \esp{de nada sirve que} + subjonctif : \esp{De nada sirve que te quejes, la decisión está tomada.} = Tu as beau te plaindre, la décision est prise. Cette tournure, très idiomatique, est excellente au thème comme au commentaire.
\end{exemplebox}

\courssub{\esp{Aunque} + indicatif ou subjonctif : le piège du « bien que »}

\begin{propriete}[Les deux valeurs de \esp{aunque}]
\begin{itemize}
\item \esp{Aunque} + \textbf{indicatif} : concession portant sur un \textbf{fait avéré, certain}. \esp{Aunque estudia mucho, no aprueba.} = Bien qu'il étudie beaucoup (c'est un fait), il ne réussit pas.
\item \esp{Aunque} + \textbf{subjonctif} : \textbf{hypothèse}, fait non réalisé ou futur. \esp{Aunque estudie mucho, quizá no apruebe.} = Même s'il étudie beaucoup, il ne réussira peut-être pas.
\end{itemize}
\end{propriete}

\begin{attention}[« Avoir beau » n'est pas « bien que »]
Le francophone traduit souvent « Il a beau étudier, il échoue » par \esp{Aunque estudie, fracasa}. C'est une \textbf{double erreur} : \esp{aunque} + subjonctif signifie « même s'il étudiait » (hypothèse), et non un effort vain constaté. Pour « avoir beau », choisissez \esp{por más que} + subjonctif, qui contient l'idée d'insistance. Réservez \esp{aunque} + indicatif à « bien que » + fait certain : \esp{Aunque estudia mucho, fracasa.} = Bien qu'il étudie beaucoup, il échoue.
\end{attention}

Comparons les trois phrases du programme :

\begin{center}
\begin{tabularx}{\linewidth}{|X|X|X|}
\hline
\rowcolor{popblueL} Français & Español & Valeur \\
\hline
Il a beau étudier, il échoue. & \esp{Por más que estudie, fracasa.} & effort vain, insistance \\
\hline
Bien qu'il étudie, il échoue. & \esp{Aunque estudia, fracasa.} & fait avéré, constat \\
\hline
Même s'il étudiait, il échouerait. & \esp{Aunque estudiara, fracasaría.} & hypothèse \\
\hline
\end{tabularx}
\end{center}

\courssub{La tournure soutenue : \esp{por} + adjectif/adverbe + \esp{que} + subjonctif}

\begin{propriete}[\esp{Por\ldots que}]
Pour insister sur le \textbf{degré} d'une qualité ou d'une manière, l'espagnol emploie \esp{por} + adjectif ou adverbe + \esp{que} + subjonctif : \esp{por rápido que corra} = il a beau courir vite / si vite qu'il coure. Registre soutenu, très apprécié au commentaire de texte et à l'écrit du Bac. Avec un nom, on ajoute \esp{mucho} : \esp{por mucho dinero que tenga}.
\end{propriete}

\begin{exemplebox}[Exemples]
\esp{Por inteligente que sea, no lo sabe todo.} = Il a beau être intelligent, il ne sait pas tout.\par
\esp{Por mucho dinero que tenga, no es feliz.} = Il a beau avoir de l'argent, il n'est pas heureux.\par
\esp{Por temprano que te levantes, el bendskin ya habrá pasado.} = Tu as beau te lever tôt, la moto-taxi sera déjà passée.
\end{exemplebox}

\courssub{Cas particulier : « avoir beau jeu de »}

\begin{propriete}[« Avoir beau jeu de » + infinitif]
Cette locution (« il est facile de\ldots », souvent ironique) se traduit par \esp{tener fácil} + infinitif ou \esp{ser fácil} + infinitif : \esp{Los críticos tienen fácil juzgar desde fuera.} = Les critiques ont beau jeu de juger de l'extérieur. \esp{Es fácil hablar cuando no se arriesga nada.} = On a beau jeu de parler quand on ne risque rien.
\end{propriete}

\begin{attention}[Sujet animé obligatoire]
« Avoir beau » implique toujours un \textbf{sujet animé qui tente quelque chose}. Il n'existe pas d'équivalent impersonnel : on ne peut pas traduire mot à mot « il a beau pleuvoir ». On animera le sujet avec le subjonctif — \esp{por más que llueva, saldremos} (« il a beau pleuvoir, nous sortirons ») — ou l'on choisira \esp{aunque} : \esp{aunque haga frío, iremos al estadio} (« il a beau faire froid, nous irons au stade »). Jamais d'indicatif après \esp{por más que} dans l'espagnol standard.
\end{attention}

\begin{popfigure}
\begin{tikzpicture}
\matrix (m) [matrix of nodes, nodes={draw,rounded corners,minimum width=4.6cm,minimum height=1.6cm,align=center,font=\small}, row sep=0.5cm, column sep=0.8cm] {
|[fill=popblueL]| \esp{por más que} + subj.\\ effort vain, choix par défaut & |[fill=popgreenL]| \esp{aunque} + indicatif\\ fait certain (« bien que ») \\
|[fill=popgoldL]| \esp{por mucho que} + subj.\\ insistance sur la quantité & |[fill=poporangeL]| \esp{aunque} + subjonctif\\ hypothèse, futur \\
|[fill=poppinkL]| \esp{por} + adj/adv + \esp{que} + subj.\\ registre soutenu, degré & |[fill=poppurpleL]| \esp{tener fácil / ser fácil}\\ « avoir beau jeu de » \\
};
\end{tikzpicture}
\legende{Carte mentale des six traductions possibles de la concession autour de « avoir beau ».}
\end{popfigure}

\begin{methode}[Traduire « avoir beau » en trois questions]
\begin{enumerate}
\item \textbf{Est-ce un effort vain d'un sujet animé ?} → \esp{por más que} / \esp{por mucho que} + subjonctif. C'est le choix par défaut.
\item \textbf{Veut-on insister sur une qualité ou une manière ?} → \esp{por} + adjectif/adverbe + \esp{que} + subjonctif (soutenu).
\item \textbf{Le français dit-il « bien que » + fait certain ?} → \esp{aunque} + indicatif. Attention : ce n'est plus « avoir beau ».
\end{enumerate}
Vérifiez toujours le mode : après \esp{por más/mucho que}, jamais d'indicatif en espagnol standard.
\end{methode}

\begin{savaistu}[Variante latino-américaine]
En Amérique latine, dans l'oral familier, on entend parfois \esp{por más que} + \textbf{indicatif} pour insister sur la réalité du fait : \esp{Por más que busqué, no lo encontré} (« j'ai eu beau chercher »). Cette variante régionale, courante au Mexique ou en Argentine, est à reconnaître à la compréhension, mais à l'examen on exige la norme : subjonctif obligatoire.
\end{savaistu}

\begin{exoresolu}[Thème guidé]
Traduisez : « Les élèves ont beau réviser toute la nuit, ils oublient tout le jour de l'examen. »\par
\tcblower
\textbf{Étape 1} : sujet animé (\emph{les élèves}) + effort vain (\emph{oublient tout}) → structure \esp{por más que} + subjonctif.\par
\textbf{Étape 2} : subjonctif présent de \esp{repasar} : \esp{repasen}.\par
\textbf{Étape 3} : \esp{Por más que repasen toda la noche, los alumnos lo olvidan todo el día del examen.}\par
\textbf{Vérification} : subjonctif après \esp{por más que} ; sujet animé ; résultat contraire à l'indicatif. Traduction correcte.
\end{exoresolu}

\begin{exercice}[À vous de traduire]
\begin{enumerate}
\item Elle a beau travailler, elle n'arrive pas à finir à temps.
\item Il a beau être riche, il n'est pas généreux. (deux versions possibles)
\item Bien qu'il pleuve, le match aura lieu à Yaoundé.
\item Les journalistes ont beau jeu de critiquer le gouvernement.
\end{enumerate}
\end{exercice}

\begin{corrige}[Exercice « À vous de traduire »]
\begin{enumerate}
\item \esp{Por más que trabaje, no consigue terminar a tiempo.} — subjonctif obligatoire après \esp{por más que}.
\item \esp{Por rico que sea, no es generoso.} ou \esp{Por mucho dinero que tenga, no es generoso.}
\item \esp{Aunque llueve, el partido tendrá lugar en Yaundé.} — fait avéré : indicatif.
\item \esp{Los periodistas tienen fácil criticar al gobierno.}
\end{enumerate}
\end{corrige}

\begin{aretenir}[L'essentiel de la section]
« Avoir beau » = \esp{por más que} / \esp{por mucho que} + \textbf{subjonctif} (effort vain, sujet animé). « Bien que » + fait certain = \esp{aunque} + \textbf{indicatif}. Registre soutenu : \esp{por} + adjectif + \esp{que} + subjonctif. « Avoir beau jeu de » = \esp{tener fácil} + infinitif. Ne confondez jamais « avoir beau » avec \esp{aunque} + subjonctif, qui exprime l'hypothèse.
\end{aretenir}

\coursec{Traduire « demander » : cuatro verbos, cuatro matices}

Dans la section précédente, nous avons vu comment rendre la concession de « avoir beau » par \esp{por mucho que}, \esp{aunque} et leurs constructions. Nous abordons maintenant un autre verbe français apparemment simple : « demander ». Simple en français, mais redoutable en espagnol, car là où le français n'a qu'un seul verbe, l'espagnol en possède quatre : \esp{pedir}, \esp{preguntar}, \esp{rogar} et \esp{solicitar}. Choisir le bon verbe, c'est déjà traduire à moitié.

\courssub{Observation : un dialogue au commissariat de Yaoundé}

Lisons d'abord ce dialogue original. Repérez les quatre verbes mis en gras et demandez-vous, pour chacun, ce que le locuteur attend en retour.

\begin{textebox}[En la comisaría]
— Buenos días, señor comisario. Vengo a \textbf{solicitar} un certificado de residencia para mi expediente de beca.\\
— Muy bien. Rellene este formulario y \textbf{pida} también una fotocopia de su cédula en la ventanilla dos.\\
— De acuerdo. Perdone, ¿puedo \textbf{preguntarle} cuánto tiempo tarda el trámite?\\
— Unos diez días. ¿Algo más?\\
— Sí, se lo \textbf{ruego}: mi hermano está enfermo en el hospital y necesito el documento antes del viernes. Le \textbf{ruego} que haga lo posible.\\
— Haré lo que pueda, pero no le prometo nada.\\
— Se lo agradezco mucho. Ah, y otra cosa: mi vecino quiere \textbf{pedirle} un favor. Me \textbf{preguntó} si usted podría recibirlo mañana.\\
— Que pase mañana a las diez y veremos.
\end{textebox}

\textbf{Glossaire :} \esp{comisaría} = commissariat ; \esp{certificado de residencia} = certificat de résidence ; \esp{expediente de beca} = dossier de bourse ; \esp{rellenar un formulario} = remplir un formulaire ; \esp{ventanilla} = guichet ; \esp{cédula} = carte d'identité ; \esp{trámite} = démarche administrative ; \esp{tardar} = prendre (du temps) ; \esp{recibirlo} = le recevoir.

\textbf{Questions de compréhension :}
\begin{enumerate}
\item Quels documents le visiteur doit-il fournir ?
\item Pourquoi demande-t-il que le dossier aille plus vite ?
\item Quelle est la différence de situation entre \esp{solicitar un certificado} et \esp{rogar que haga lo posible} ?
\end{enumerate}

On observe que chaque verbe correspond à une intention différente : une démarche officielle (\esp{solicitar}), un objet ou un service attendu (\esp{pedir}), une information (\esp{preguntar}), une supplication émotive (\esp{rogar}). C'est exactement la logique qu'il faut acquérir pour le thème comme pour la version.

\courssub{Pedir : demander pour recevoir}

\begin{definition}[Pedir]
\esp{Pedir} signifie demander quelque chose que l'on attend de \cle{recevoir} : un objet concret, un service, une action, une permission. C'est un verbe \cle{transitif direct} : on demande \emph{quelque chose} \emph{à quelqu'un}.
\end{definition}

\begin{propriete}[Construction et conjugaison de pedir]
La construction est : \esp{pedir algo a alguien} (demander quelque chose à quelqu'un). Attention : la personne est introduite par \esp{a}, jamais par \esp{de}.\\
\esp{Pedir} est un verbe à changement radical e $\rightarrow$ i : \esp{pido, pides, pide, pedimos, pedís, piden} ; prétérit : \esp{pedí, pediste, pidió, pedimos, pedisteis, pidieron} ; gérondif : \esp{pidiendo}.\\
Suivi d'une subordonnée d'action, il exige le subjonctif : \esp{Te pido que vengas.} (Je te demande de venir.) On ne peut pas dire \esp{te pido venir} quand les sujets sont différents : l'infinitif n'est possible que si le sujet est le même (\esp{Pido salir.} = Je demande à sortir, moi-même).
\end{propriete}

\begin{exemplebox}[Pedir en contexte]
\esp{Le pidió un favor a su jefe.} = Il a demandé un service à son patron.\\
\esp{Los alumnos piden permiso para salir.} = Les élèves demandent la permission de sortir.\\
\esp{Pedí la cuenta al camarero.} = J'ai demandé l'addition au serveur.\\
\esp{Mi madre me pidió que comprara pan en el mercado de Mokolo.} = Ma mère m'a demandé d'acheter du pain au marché de Mokolo.\\
\esp{El futbolista pidió el balón.} = Le footballeur a demandé le ballon.
\end{exemplebox}

\courssub{Preguntar : demander pour savoir}

\begin{definition}[Preguntar]
\esp{Preguntar} signifie demander une \cle{information}, poser une question : on attend de \emph{recevoir des mots}, une réponse verbale. Construction : \esp{preguntar algo a alguien} ou \esp{preguntar por algo/alguien} (demander des nouvelles de, s'enquérir de).
\end{definition}

\begin{propriete}[Preguntar et l'interrogation indirecte]
Après \esp{preguntar}, l'interrogation indirecte garde l'\cle{indicatif} (c'est une question rapportée, pas une volonté) et le mot interrogatif prend un accent :\\
\esp{Me preguntó si estaba cansado.} = Il m'a demandé si j'étais fatigué.\\
\esp{Le pregunté dónde vivía.} = Je lui ai demandé où il habitait.\\
Ne confondez pas : \esp{pedir que + subjonctif} (je veux une action) mais \esp{preguntar si/cuándo/dónde + indicatif} (je veux une information).
\end{propriete}

\begin{exemplebox}[Preguntar en contexte]
\esp{Le pregunté la hora.} = Je lui ai demandé l'heure.\\
\esp{El profesor preguntó quién había roto la ventana.} = Le professeur a demandé qui avait cassé la fenêtre.\\
\esp{Preguntó por ti esta mañana.} = Il a demandé de tes nouvelles ce matin.\\
\esp{Pregunta el camino a ese taxista.} = Demande le chemin à ce chauffeur de taxi.
\end{exemplebox}

\begin{methode}[La question-clé pour choisir entre pedir et preguntar]
Avant de traduire « demander », pose-toi cette question unique :
\begin{enumerate}
\item Est-ce que j'attends un \textbf{objet, un service ou une action} ? $\rightarrow$ \esp{pedir}.
\item Est-ce que j'attends une \textbf{réponse verbale, une information} ? $\rightarrow$ \esp{preguntar}.
\end{enumerate}
Test : « Il demande le prix. » J'attends une information (un chiffre dit à voix haute) $\rightarrow$ \esp{Pregunta el precio.} En revanche : « Il demande une augmentation. » J'attends une décision, un avantage $\rightarrow$ \esp{Pide un aumento.}
\end{methode}

\begin{attention}[Pièges des francophones avec pedir et preguntar]
\begin{itemize}
\item « Demander en mariage » se dit \esp{pedir la mano} (locution figée), jamais \esp{preguntar} : \esp{Le pidió la mano en diciembre.} = Il l'a demandée en mariage en décembre.
\item « Demander son chemin » se dit \esp{preguntar el camino} (on attend une indication orale), pas \esp{pedir el camino}.
\item Ne jamais dire \esp{pedir una pregunta} au sens de « poser une question » : on dit \esp{hacer una pregunta}. \esp{Pedir la palabra} = demander la parole (pour la prendre).
\item Faux ami : \esp{demandar} existe en espagnol mais signifie « poursuivre en justice » ! \esp{Demandaron a la empresa.} = Ils ont poursuivi l'entreprise en justice. Ne l'utilisez jamais pour traduire le « demander » ordinaire.
\end{itemize}
\end{attention}

\courssub{Rogar : supplier, implorer}

\begin{definition}[Rogar]
\esp{Rogar} exprime la demande chargée d'\cle{émotion} : supplier, implorer, prier instamment. Registre soutenu ou affectif. C'est un verbe à changement radical o $\rightarrow$ ue : \esp{ruego, ruegas, ruega, rogamos, rogáis, ruegan}. Attention au prétérit, avec changement orthographique g $\rightarrow$ gu devant e : \esp{rogué, rogaste, rogó, rogamos, rogasteis, rogaron}.
\end{definition}

\begin{propriete}[Rogar et le subjonctif]
\esp{Rogar} exprime une volonté, une influence sur autrui : il déclenche le \cle{subjonctif} dans la subordonnée introduite par \esp{que}. Construction : \esp{rogar a alguien que + subjuntivo}.\\
\esp{Te ruego que me escuches.} = Je t'en supplie, écoute-moi.\\
\esp{Les rogamos que no hagan ruido.} = Nous vous prions de ne pas faire de bruit. (formule administrative de politesse, très fréquente sur les panneaux : \esp{Se ruega no fumar} = Prière de ne pas fumer.)
\end{propriete}

\begin{exemplebox}[Rogar en contexte]
\esp{El enfermo rogó al médico que lo curara.} = Le malade a supplié le médecin de le guérir.\\
\esp{Te ruego un favor: acompáñame al hospital.} = Je te demande instamment un service : accompagne-moi à l'hôpital.\\
\esp{La madre rogaba por la vida de su hijo.} = La mère suppliait pour la vie de son fils.
\end{exemplebox}

\courssub{Solicitar : la demande formelle et administrative}

\begin{definition}[Solicitar]
\esp{Solicitar} signifie demander de manière \cle{formelle}, le plus souvent par écrit, dans le cadre d'une procédure administrative ou professionnelle : un emploi, un visa, une bourse, un document. Registre formel. Attention : \esp{solicitar} ne signifie jamais « solliciter » au sens de « importuner ».
\end{definition}

\begin{exemplebox}[Solicitar en contexte]
\esp{Solicité el visado por internet.} = J'ai demandé le visa en ligne.\\
\esp{Ha solicitado una beca para estudiar en Madrid.} = Il a demandé une bourse pour étudier à Madrid.\\
\esp{Solicitamos que se nos envíe el expediente.} = Nous demandons que le dossier nous soit envoyé. (subjonctif après \esp{solicitar que}, comme après \esp{pedir que})
\end{exemplebox}

\begin{propriete}[Le passif impersonnel avec solicitar]
\esp{Solicitar} apparaît très souvent dans la construction passive réflexive \esp{se solicita}, typique des annonces :\\
\esp{Se solicita personal.} = On demande du personnel.\\
\esp{Se solicitan profesores de español.} = On recherche des professeurs d'espagnol.\\
C'est la tournure des petites annonces et des avis officiels ; le verbe s'accorde avec le nom qui suit (\esp{se solicita} / \esp{se solicitan}).
\end{propriete}

\begin{savaistu}[Pedir la cuenta : Espagne et Amérique]
En Espagne, au restaurant, on dit \esp{¿Me trae la cuenta, por favor?} ou simplement \esp{La cuenta, por favor}. En Amérique latine, on entend souvent la formule sans verbe : \esp{La cuenta, por favor}. Utiliser \esp{pedir la cuenta} est correct partout, mais les Hispanophones économisent volontiers le verbe quand le contexte est clair. Autre détail : en Guinée équatoriale, seul pays hispanophone d'Afrique subsaharienne, les démarches administratives utilisent massivement \esp{solicitar} et \esp{se ruega}, hérités de la tradition administrative espagnole.
\end{savaistu}

\courssub{Synthèse visuelle et tableau récapitulatif}

\begin{popfigure}
\begin{tikzpicture}[node distance=1.2cm]
\node[draw,rounded corners,fill=popblueL,align=center] (pedir) {PEDIR\\Objet / Service\\« Je veux avoir »};
\node[draw,rounded corners,fill=popgreenL,right=2cm of pedir,align=center] (preguntar) {PREGUNTAR\\Information\\« Je veux savoir »};
\node[draw,rounded corners,fill=popgoldL,below=1.5cm of pedir,align=center] (rogar) {ROGAR\\Supplication\\« Je supplie »};
\node[draw,rounded corners,fill=poporangeL,right=2cm of rogar,align=center] (solicitar) {SOLICITAR\\Formel / Admin.\\« Je fais une demande »};
\draw[<->,dashed,popdark] (pedir) -- node[above,align=center] {\small Objet vs Info} (preguntar);
\draw[<->,dashed,popdark] (rogar) -- node[above,align=center] {\small Émotion vs Procédure} (solicitar);
\draw[<->,dashed,popmuted] (pedir) -- node[left,align=center] {\small Neutre vs Émotif} (rogar);
\draw[<->,dashed,popmuted] (preguntar) -- node[right,align=center] {\small Oral vs Écrit} (solicitar);
\end{tikzpicture}
\legende{Carte mentale : les quatre verbes espagnols pour « demander » et leurs oppositions.}
\end{popfigure}

\begin{center}
\begin{tabularx}{\linewidth}{|l|X|X|X|}
\hline
\rowcolor{popblueL} \textbf{Verbe} & \textbf{On attend...} & \textbf{Construction} & \textbf{Exemple} \\
\hline
\esp{pedir} & un objet, un service, une action & \esp{pedir algo a alguien} ; \esp{pedir que + subj.} & \esp{Pidió un favor.} (Il a demandé un service.) \\
\hline
\esp{preguntar} & une information, une réponse & \esp{preguntar algo a alguien} ; \esp{preguntar por} ; \esp{preguntar si + ind.} & \esp{Preguntó la hora.} (Il a demandé l'heure.) \\
\hline
\esp{rogar} & une grâce, avec émotion & \esp{rogar a alguien que + subj.} & \esp{Te ruego que vengas.} (Je te supplie de venir.) \\
\hline
\esp{solicitar} & un droit, par procédure formelle & \esp{solicitar algo} ; \esp{se solicita...} & \esp{Solicitó el pasaporte.} (Il a demandé le passeport.) \\
\hline
\end{tabularx}
\end{center}

Contrastes à mémoriser : \esp{pedir perdón} (demander pardon : on attend le pardon, une grâce) mais \esp{preguntar la hora} (demander l'heure : on attend une information) ; \esp{rogar un favor} (supplier pour obtenir un service, ton ému) mais \esp{solicitar un visado} (demande officielle d'un visa, ton administratif).

\courssub{Exercice résolu et entraînement}

\begin{exoresolu}[Choix du verbe et traduction]
Traduisez en espagnol : 1. Elle a demandé un verre d'eau au serveur. 2. Le journaliste a demandé au ministre quand commencerait la réunion. 3. Je vous en supplie, aidez ma famille. 4. On demande des chauffeurs expérimentés. 5. Il m'a demandé de lui prêter mon vélo.
\tcblower
\textbf{Solution détaillée :}
\begin{enumerate}
\item Objet attendu $\rightarrow$ \esp{pedir} : \esp{Le pidió un vaso de agua al camarero.} (attention : \esp{al camarero}, la personne est introduite par \esp{a}).
\item Information attendue $\rightarrow$ \esp{preguntar} + interrogation indirecte : \esp{El periodista le preguntó al ministro cuándo empezaría la reunión.} (accent sur \esp{cuándo} en interrogation indirecte ; conditionnel \esp{empezaría} pour la concordance des temps après un principal au passé).
\item Supplication $\rightarrow$ \esp{rogar que + subj.} : \esp{Se lo ruego, ayude a mi familia.} ou \esp{Les ruego que ayuden a mi familia.}
\item Annonce formelle $\rightarrow$ passif réflexif : \esp{Se solicitan conductores experimentados.} (accord au pluriel avec \esp{conductores}).
\item Action attendue $\rightarrow$ \esp{pedir que + subj.} : \esp{Me pidió que le prestara mi bicicleta.} (imparfait du subjonctif \esp{prestara} après une principale au passé).
\end{enumerate}
\end{exoresolu}

\begin{exercice}[À vous de traduire]
Traduisez : 1. Demande ton chemin au policier. 2. Les élèves ont demandé la parole. 3. Il a demandé sa fille en mariage. 4. Nous demandons que le directeur nous reçoive. 5. Le malade suppliait le docteur de le sauver. 6. J'ai demandé un duplicata de mon acte de naissance à la mairie de Douala.
\end{exercice}

\begin{corrige}[Exercice : les quatre verbes]
1. \esp{Pregunta el camino al policía.} (information $\rightarrow$ \esp{preguntar}).\\
2. \esp{Los alumnos pidieron la palabra.} (on attend de la prendre $\rightarrow$ \esp{pedir}).\\
3. \esp{Pidió la mano de su hija.} (locution figée \esp{pedir la mano}).\\
4. \esp{Solicitamos que el director nos reciba.} (demande formelle + subjonctif \esp{reciba}).\\
5. \esp{El enfermo rogaba al médico que lo salvara.} (\esp{rogar} + imparfait du subjonctif).\\
6. \esp{Solicité un duplicado de mi acta de nacimiento en el ayuntamiento de Douala.} (procédure administrative $\rightarrow$ \esp{solicitar}).
\end{corrige}

\begin{aretenir}[L'essentiel de la section]
« Demander » n'a pas de traduction unique : \esp{pedir} pour recevoir (objet, service, action), \esp{preguntar} pour savoir (information), \esp{rogar} pour supplier (émotion + subjonctif), \esp{solicitar} pour une démarche formelle (administration, annonces \esp{se solicita}). La question-clé : « Qu'est-ce que j'attends en retour ? » Locutions figées à connaître par cœur : \esp{pedir la mano}, \esp{pedir perdón}, \esp{pedir la palabra}, \esp{preguntar la hora}, \esp{preguntar el camino}, \esp{hacer una pregunta}. Et méfiez-vous du faux ami \esp{demandar} : « poursuivre en justice ».
\end{aretenir}

\coursec{Traduire « depuis » : repères temporels et aspect}

Après avoir appris à nuancer la demande avec \esp{pedir}, \esp{preguntar}, \esp{rogar} et \esp{solicitar}, nous abordons maintenant une petite préposition française qui cache, en espagnol, tout un système : \cle{depuis}. Selon que « depuis » renvoie à une \emph{date}, à une \emph{durée} ou à une \emph{proposition entière}, l'espagnol choisit des structures différentes : \esp{desde}, \esp{desde hace}, \esp{hace\ldots que} ou \esp{desde que}. C'est l'une des sources d'erreurs les plus fréquentes au baccalauréat : observons d'abord un texte.

\courssub{Observation : un texte pour voir « depuis » en action}

\begin{textebox}[Un periodista de Yaoundé]
Trabajo en esta redacción desde 2018. Desde que llegué, he cubierto elecciones, partidos de fútbol y festivales de música. Hace seis años que escribo cada semana una crónica sobre la vida de los barrios: los mercados, los bendskins, las escuelas. Vivo en este barrio desde hace diez años y, desde que nacieron mis hijos, miro la ciudad con otros ojos. Desde el primer día, comprendí que informar es también construir la paz: desde que publiqué mi primer artículo sobre el diálogo entre comunidades, recibo cartas de lectores de todo el país. Mi jefe me dice: « Desde que tú estás aquí, el periódico ha cambiado ». Y yo respondo que seguiré escribiendo desde que amanezca hasta que anochezca, porque esta profesión es mi vida.
\end{textebox}

\textbf{Glossaire :} \esp{redacción} = rédaction (d'un journal) ; \esp{cubrir} = couvrir (un événement) ; \esp{crónica} = chronique ; \esp{barrio} = quartier ; \esp{amanecer} = (au lever du jour) ; \esp{anochecer} = à la tombée de la nuit ; \esp{jefe} = chef, patron.

\textbf{Questions de compréhension :}
\begin{enumerate}
\item Repérez toutes les expressions qui traduisent « depuis » ou « il y a\ldots que ».
\item Pour chacune, dites si le repère est une \emph{date}, une \emph{durée} ou une \emph{proposition avec verbe}.
\item Quels temps de verbe suivent \esp{desde que} dans ce texte ?
\end{enumerate}

\textbf{Observations.} On remarque trois grandes configurations :
\begin{itemize}
\item \esp{desde 2018}, \esp{desde el primer día} : \esp{desde} + un \emph{nom ou une date} (point de départ).
\item \esp{desde hace diez años}, \esp{hace seis años que escribo} : une \emph{durée} écoulée.
\item \esp{desde que llegué}, \esp{desde que nacieron mis hijos} : \esp{desde que} + une \emph{proposition verbale}.
\end{itemize}

\courssub{Desde + point de départ : date ou événement}

\begin{propriete}[Desde + nom / date]
\esp{Desde} est une \cle{préposition} suivie d'un nom, d'une date ou d'un événement. Elle marque le \emph{point de départ} d'une action ou d'un état. Si cette action \emph{continue encore au moment où l'on parle}, le verbe principal se met à l'\cle{indicatif présent} — là où le français utilise aussi le présent : \esp{Vivo en Douala desde 2015.} = « J'habite Douala depuis 2015. »
\end{propriete}

\begin{exemplebox}[Desde + point de départ]
\begin{itemize}
\item \esp{Trabajo aquí desde enero.} = « Je travaille ici depuis janvier. »
\item \esp{Estudia español desde la clase de seconde.} = « Il étudie l'espagnol depuis la seconde. »
\item \esp{Desde la independencia, el país ha cambiado mucho.} = « Depuis l'indépendance, le pays a beaucoup changé. »
\item \esp{No la veo desde el lunes.} = « Je ne l'ai pas vue depuis lundi. »
\item \esp{Desde pequeño, sueña con ser periodista.} = « Depuis tout petit, il rêve d'être journaliste. »
\end{itemize}
\end{exemplebox}

\begin{attention}[Présent, pas passé composé !]
Le francophone traduit souvent « J'habite ici depuis 2015 » par un passé en espagnol. Erreur ! Si l'action continue, l'espagnol exige le \emph{présent} : \esp{Vivo aquí desde 2015} (et non \esp{he vivido}). En revanche, si l'action est terminée, on utilise le passé : \esp{Viví allí desde 2010 hasta 2015.} = « J'y ai vécu de 2010 à 2015. »
\end{attention}

\courssub{Desde hace + durée / Hace + durée + que : le focus sur la durée}

Quand « depuis » introduit une \emph{durée} (« depuis dix ans »), l'espagnol ne peut pas dire \esp{desde diez años}. Il faut ajouter \esp{hace} :

\begin{propriete}[Les deux constructions de la durée]
Pour exprimer « depuis + durée » avec une action encore en cours :
\begin{itemize}
\item \esp{desde hace + durée} : \esp{Vivo aquí desde hace diez años.}
\item \esp{hace + durée + que + présent} : \esp{Hace diez años que vivo aquí.}
\end{itemize}
Les deux phrases signifient : « J'habite ici depuis dix ans. » La seconde construction met l'accent sur la durée écoulée.
\end{propriete}

\begin{exemplebox}[Durée : les deux formes]
\begin{itemize}
\item \esp{Hace dos horas que espero el autobús.} = « J'attends le bus depuis deux heures. » (ou : \esp{Espero el autobús desde hace dos horas.})
\item \esp{Hace mucho tiempo que no vemos un partido así.} = « Il y a longtemps que nous n'avons pas vu un tel match. »
\item \esp{Llueve desde hace tres días.} = « Il pleut depuis trois jours. »
\item \esp{Hace un mes que estudian el subjuntivo.} = « Ils étudient le subjonctif depuis un mois. »
\end{itemize}
\end{exemplebox}

\begin{attention}[Piège majeur : *desde diez años]
On ne dit JAMAIS \esp{desde diez años}. Devant une durée, deux solutions seulement : \esp{desde hace diez años} ou \esp{hace diez años que}. Autre piège : « Il y a dix ans qu'il habite ici » ne se traduit pas par \esp{desde} seul, mais par \esp{Hace diez años que vive aquí}. Le mot \esp{hace} est invariable : jamais de \esp{hacen} même devant un pluriel (\esp{hace diez años}).
\end{attention}

\courssub{Desde que + verbe : indicatif ou subjonctif}

Quand « depuis » est suivi d'un \emph{verbe conjugué} (« depuis que je suis arrivé »), le français utilise « depuis que » ; l'espagnol utilise \esp{desde que}, qui est alors une \cle{conjonction}.

\begin{propriete}[Desde que : le choix du mode]
\begin{itemize}
\item \esp{Desde que + indicatif} : le fait est \emph{certain, réel, passé}. Le verbe est souvent au \emph{pretérito indefinido} ou au \emph{perfecto} : \esp{Desde que llegué, no he salido.} = « Depuis que je suis arrivé, je ne suis pas sorti. »
\item \esp{Desde que + subjonctif} : le fait est \emph{futur ou hypothétique} (complément niveau Bac) : \esp{Te esperaré desde que llegues.} = « Je t'attendrai dès que tu arriveras / à partir de ton arrivée. »
\end{itemize}
\end{propriete}

\begin{exemplebox}[Desde que + indicatif (fait réel)]
\begin{itemize}
\item \esp{Desde que nací, mi familia vive en esta casa.} = « Depuis que je suis né, ma famille vit dans cette maison. »
\item \esp{Desde que llegué, no he salido.} = « Depuis que je suis arrivé, je ne suis pas sorti. »
\item \esp{Desde que empezó a estudiar español, comprende mejor las noticias.} = « Depuis qu'il a commencé à étudier l'espagnol, il comprend mieux les informations. »
\item \esp{Desde que se conocieron, son inseparables.} = « Depuis qu'ils se sont connus, ils sont inséparables. »
\end{itemize}
\end{exemplebox}

\begin{exemplebox}[Desde que + subjonctif (futur / hypothétique)]
\begin{itemize}
\item \esp{Te esperaré desde que llegues.} = « Je t'attendrai dès ton arrivée. »
\item \esp{Desde que tengas tiempo, ven a verme.} = « Dès que tu auras le temps, viens me voir. »
\item \esp{Contaremos la historia desde que empiece la película.} = « Nous raconterons l'histoire à partir du moment où le film commencera. »
\end{itemize}
\end{exemplebox}

\begin{attention}[Depuis / depuis que : préposition ou conjonction ?]
Comme en français, il faut distinguer :
\begin{itemize}
\item \esp{desde} + \emph{nom} : \esp{desde 2010}, \esp{desde el lunes} (« depuis 2010 », « depuis lundi ») ;
\item \esp{desde que} + \emph{verbe conjugué} : \esp{desde que llegué} (« depuis que je suis arrivé »).
\end{itemize}
Ne dites jamais \esp{desde llegué} : devant un verbe, le \esp{que} est obligatoire.
\end{attention}

\courssub{Schéma récapitulatif : la ligne du temps}

\begin{popfigure}
\begin{tikzpicture}
\draw[->, thick, popdark] (0,0) -- (10,0) node[right] {Temps};
\foreach \x/\label in {0/Passé, 5/Desde, 10/Présent} \draw[popdark] (\x,0.1) -- (\x,-0.1) node[below] {\label};
\node[above=0.5cm of {(5,0)}, align=center, text=popblue] (desde) {Desde 2010 / Desde que llegué};
\node[above=1.2cm of {(5,0)}, align=center, text=poppurple] (dur) {Desde hace 10 años / Hace 10 años que\ldots};
\draw[<->, thick, poporange] (0,-1) -- node[below, text=poporange] {Durée écoulée} (10,-1);
\fill[popgreen] (5,0) circle (0.08);
\fill[popink] (10,0) circle (0.08);
\end{tikzpicture}
\legende{Le point de départ (\esp{desde}) est dans le passé ; l'action se prolonge jusqu'au présent, d'où l'indicatif présent pour une action encore en cours.}
\end{popfigure}

\courssub{Tableau récapitulatif français / espagnol}

\begin{center}
\begin{tabularx}{\linewidth}{|X|X|X|}
\hline
\rowcolor{popblueL} Français & Español & Remarque \\
\hline
depuis 2015 (date) & desde 2015 & action continue : présent \\
\hline
depuis lundi & desde el lunes & article obligatoire \\
\hline
depuis dix ans (durée) & desde hace diez años & jamais *desde diez años \\
\hline
il y a dix ans qu'il habite ici & hace diez años que vive aquí & \esp{hace} invariable \\
\hline
depuis que je suis arrivé & desde que llegué & indicatif (fait réel) \\
\hline
dès que tu arriveras & desde que llegues & subjonctif (futur) \\
\hline
depuis quand ? & ¿desde cuándo? & accent sur \esp{cuándo} \\
\hline
\end{tabularx}
\end{center}

\begin{methode}[Traduire « depuis » en quatre étapes]
\begin{enumerate}
\item \textbf{Identifier le repère} : est-ce une date ou un événement nominal (\esp{desde}), une durée (\esp{desde hace} / \esp{hace\ldots que}) ou une proposition avec verbe (\esp{desde que}) ?
\item \textbf{Vérifier si l'action est encore en cours} : si oui, verbe principal au présent de l'indicatif (\esp{vivo}, \esp{trabajo}, \esp{espero}).
\item \textbf{Si « depuis que » + verbe} : fait passé et certain $\rightarrow$ indicatif (\esp{desde que llegué}) ; fait futur ou hypothétique $\rightarrow$ subjonctif (\esp{desde que vengas}).
\item \textbf{Contrôler les pièges} : \esp{hace} invariable, \esp{que} obligatoire devant un verbe, accent dans la question \esp{¿desde cuándo?}.
\end{enumerate}
\end{methode}

\begin{exoresolu}[Thème guidé]
Traduisez en espagnol : « Il y a cinq ans que ma sœur apprend l'espagnol. Depuis qu'elle est arrivée à Yaoundé, elle n'a pas manqué un seul cours. Elle habite près du lycée depuis 2020. »
\tcblower
\textbf{Solution détaillée.}
\begin{enumerate}
\item « Il y a cinq ans que\ldots » : durée + action en cours $\rightarrow$ \esp{hace\ldots que} + présent : \esp{Hace cinco años que mi hermana aprende español.} (possible aussi : \esp{Mi hermana aprende español desde hace cinco años.})
\item « Depuis qu'elle est arrivée » : proposition verbale, fait passé certain $\rightarrow$ \esp{desde que} + indefinido ; « elle n'a pas manqué » : bilan jusqu'au présent $\rightarrow$ perfecto : \esp{Desde que llegó a Yaundé, no ha faltado a ninguna clase.}
\item « depuis 2020 » : date + action en cours $\rightarrow$ \esp{desde} + présent : \esp{Vive cerca del liceo desde 2020.}
\end{enumerate}
\textbf{Traduction complète :} \esp{Hace cinco años que mi hermana aprende español. Desde que llegó a Yaundé, no ha faltado a ninguna clase. Vive cerca del liceo desde 2020.}
\end{exoresolu}

\begin{savaistu}[¿Desde cuándo? : la norme de l'accent]
En espagnol standard, la question directe s'écrit avec un accent : \esp{¿Desde cuándo vives aquí?} (« Depuis quand habites-tu ici ? »), car \esp{cuándo} est un mot interrogatif. \textbf{La norme exige aussi l'accent dans les questions indirectes} : \esp{No sé desde cuándo vive aquí.} (et non *\esp{desde cuando}). L'omission de l'accent (\esp{desde cuando}) est un trait familier ou dialectal (Caraïbes) à éviter à l'examen.
\end{savaistu}

\begin{exercice}[Entraînement au thème]
Traduisez en espagnol :
\begin{enumerate}
\item J'attends une réponse depuis trois semaines.
\item Depuis que je suis né, ma grand-mère vit avec nous.
\item Il y a deux mois que les élèves préparent le débat.
\item Nous nous connaissons depuis la classe de première.
\item Je t'écrirai dès que tu arriveras à Douala.
\end{enumerate}
\end{exercice}

\begin{corrige}[Exercice : « depuis »]
\begin{enumerate}
\item \esp{Espero una respuesta desde hace tres semanas.} (ou : \esp{Hace tres semanas que espero una respuesta.}) — durée, action en cours.
\item \esp{Desde que nací, mi abuela vive con nosotros.} — fait passé certain : indefinido \esp{nací} ; action qui continue : présent \esp{vive}.
\item \esp{Hace dos meses que los alumnos preparan el debate.} — « il y a\ldots que » : \esp{hace\ldots que} + présent.
\item \esp{Nos conocemos desde la clase de primero.} — point de départ nominal : \esp{desde} + présent.
\item \esp{Te escribiré desde que llegues a Duala.} — fait futur : subjonctif \esp{llegues}.
\end{enumerate}
\end{corrige}

\begin{aretenir}[L'essentiel sur « depuis »]
\begin{itemize}
\item Date ou événement nominal : \esp{desde} + nom, verbe au présent si l'action continue.
\item Durée : \esp{desde hace + durée} ou \esp{hace + durée + que + présent} ; jamais \esp{desde diez años}.
\item Proposition verbale : \esp{desde que} + indicatif (passé réel) ou + subjonctif (futur, hypothétique).
\item Question : \esp{¿desde cuándo?} avec accent (aussi en question indirecte).
\end{itemize}
\end{aretenir}

\coursec{Synthèse comparative, pièges et entraînement thème/version}

Nous venons d'étudier séparément trois notions délicates pour le francophone : la concession avec \esp{por más que} et \esp{aunque}, les quatre visages de « demander » (\esp{pedir}, \esp{preguntar}, \esp{rogar}, \esp{solicitar}) et les repères temporels de « depuis » (\esp{desde}, \esp{desde hace}, \esp{hace\ldots que}, \esp{desde que}). Il est temps de rassembler tout cela dans un tableau unique, de traquer les erreurs d'interférence, puis de nous entraîner dans les deux sens de la traduction, comme au baccalauréat.

\courssub{Tableau récapitulatif des trois notions}

\begin{aretenir}[Tableau récapitulatif FR / ES à photocopier]
Ce tableau synthétise les choix de traduction, le mode verbal et la structure exigés pour chaque cas. Recopiez-le dans votre cahier de grammaire : il condense toute la leçon.
\end{aretenir}

\begin{popfigure}
\begin{tikzpicture}
\matrix (tab) [matrix of nodes, nodes={draw=popline, minimum width=3.2cm, minimum height=0.9cm, align=center, font=\small}, row sep=0.15cm, column sep=0.15cm] {
|[fill=popblueL]| FRANÇAIS & |[fill=popblueL]| ESPAGNOL & |[fill=popblueL]| MODE / STRUCTURE \\
avoir beau + inf. & por más que / por mucho que & + subjonctif présent \\
bien que (fait réel) & aunque & + indicatif \\
demander (une chose) & pedir & pedir algo a alguien \\
demander (une info) & preguntar & preguntar algo / por \\
supplier, prier & rogar & rogar que + subjonctif \\
demander (formel) & solicitar & + nom / que + subj. \\
depuis (date, moment) & desde + nom / date & présent de l'indicatif \\
depuis (durée) & desde hace / hace\ldots que & présent de l'indicatif \\
depuis que (passé) & desde que & + indicatif passé \\
depuis que (futur) & desde que & + subjonctif \\
};
\end{tikzpicture}
\legende{Synthèse des trois notions : choix lexical, mode et structure.}
\end{popfigure}

\begin{center}
\begin{tabularx}{\linewidth}{|X|X|X|}
\hline
\rowcolor{popblueL} Français & Español & Remarque \\
\hline
Il a beau essayer & Por mucho que lo intente & subjonctif obligatoire \\
\hline
Bien qu'il pleuve (c'est un fait) & Aunque llueve & indicatif : fait constaté \\
\hline
Je demande du pain & Pido pan & \esp{pedir} + chose \\
\hline
Je demande l'heure & Pregunto la hora & \esp{preguntar} + information \\
\hline
Il habite ici depuis 2010 & Vive aquí desde 2010 & présent espagnol \\
\hline
Il habite ici depuis dix ans & Vive aquí desde hace diez años & jamais \esp{*desde diez años} \\
\hline
\end{tabularx}
\end{center}

\courssub{Pièges majeurs et interférences du français}

\begin{attention}[Les cinq pièges majeurs]
\begin{enumerate}
\item Ne jamais traduire « demander » par \esp{*demandar} : ce verbe signifie « exiger, réclamer avec force, poursuivre en justice ». \esp{El sindicato demanda mejores salarios.} « Le syndicat exige de meilleurs salaires. »
\item Ne jamais dire \esp{*desde diez años} pour « depuis dix ans » : avec une durée, il faut \esp{desde hace diez años} ou \esp{hace diez años que\ldots} La structure \esp{desde} + durée seule est un calque du français.
\item « Avoir beau » ne se traduit pas littéralement : \esp{*tener bello} n'existe pas. Il faut \esp{por más que} / \esp{por mucho que} + subjonctif, ou \esp{aunque} + subjonctif si le fait est hypothétique.
\item \esp{Desde que} suivi d'un fait futur ou non accompli exige le subjonctif : \esp{Desde que llegues, empezamos.} « Dès que tu arriveras, nous commencerons. » L'indicatif est réservé au passé : \esp{Desde que llegó, todo cambió.} « Depuis qu'il est arrivé, tout a changé. »
\item Pour demander une information, on n'utilise pas \esp{pedir} : \esp{*pedir la hora} est un calque du français ; on dit \esp{preguntar la hora}, \esp{preguntar por el camino}.
\end{enumerate}
\end{attention}

\begin{exemplebox}[Discussion de nuance : « avoir beau »]
Phrase de thème : « Il a beau essayer, il n'y arrive pas. »

Deux traductions possibles :
\begin{itemize}
\item \esp{Aunque lo intente, no lo logra.} — concession hypothétique, registre courant.
\item \esp{Por más que lo intente, no lo logra.} — insistance : « quels que soient ses efforts ». C'est la traduction la plus fidèle de « avoir beau », qui exprime l'inutilité de l'effort.
\end{itemize}
La nuance : \esp{aunque} constate un obstacle ; \esp{por más que} souligne l'intensité vaine de l'action. Au bac, préférez \esp{por más/mucho que} dès que « avoir beau » apparaît.
\end{exemplebox}

\begin{exemplebox}[Version commentée]
\esp{El reportero solicita información desde hace una hora.}

Traduction : « Le journaliste demande des informations depuis une heure. »

Analyse : \esp{solicitar} = « demander » au registre formel (et non « solliciter » au sens français de « importuner ») ; \esp{desde hace una hora} = durée jusqu'au présent, donc présent en espagnol comme en français ; \esp{reportero} = « reporter, journaliste », faux ami partiel de « rapporteur ».
\end{exemplebox}

\courssub{Méthode de traduction en quatre étapes}

\begin{methode}[Traduire une phrase de thème contenant nos trois notions]
\begin{enumerate}
\item \textbf{Repérer la notion} : la phrase contient-elle « avoir beau », « demander » ou « depuis » ? Souligner le déclencheur.
\item \textbf{Choisir le verbe ou la structure} : \esp{pedir} ou \esp{preguntar} ? \esp{desde} ou \esp{desde hace} ? \esp{aunque} ou \esp{por más que} ?
\item \textbf{Accorder le mode} : subjonctif après \esp{por más que}, \esp{rogar que}, \esp{desde que} + futur ; indicatif après \esp{aunque} (fait réel) et \esp{desde que} + passé.
\item \textbf{Vérifier l'ordre des mots et la morphologie} : place de la subordonnée, accents (\esp{información}, \esp{está}), prépositions (\esp{preguntar por}, \esp{pedir algo a alguien}).
\end{enumerate}
\end{methode}

\begin{exoresolu}[Thème guidé : phrase complète]
Énoncé : Traduire en espagnol : « Elle a beau demander le chemin au vendeur, elle ne comprend pas ses explications. »
\tcblower
Étape 1 : deux notions : « avoir beau » et « demander » (un chemin = information).

Étape 2 : « avoir beau » $\rightarrow$ \esp{por más que} ; « demander le chemin » $\rightarrow$ \esp{preguntar por el camino} (information, donc \esp{preguntar}, avec \esp{por}) ; « au vendeur » $\rightarrow$ \esp{al vendedor}.

Étape 3 : subjonctif présent après \esp{por más que} : \esp{pregunte}.

Étape 4 : assemblage : \esp{Por más que le pregunte por el camino al vendedor, no entiende sus explicaciones.}

Vérification : \esp{explicaciones} sans accent (pluriel en \esp{-nes}) ; \esp{entiende} (e $\rightarrow$ ie) correct ; le pronom \esp{le} reprend \esp{al vendedor}.
\end{exoresolu}

\courssub{Entraînement au thème (FR $\rightarrow$ ES) : dix phrases progressives}

\begin{exercice}[Thème : traduisez en espagnol]
Phrases isolées :
\begin{enumerate}
\item Il a beau travailler, il échoue toujours.
\item Je demande du riz au marché de Mokolo.
\item Elle demande l'heure au chauffeur du bendskin.
\item Nous habitons à Yaoundé depuis 2015.
\item Mon père vend du cacao depuis vingt ans.
\end{enumerate}
Phrases en contexte :
\begin{enumerate}
\setcounter{enumi}{5}
\item Depuis que le journaliste est arrivé, tout le monde demande des informations.
\item Le ministre a beau parler, la foule ne l'écoute pas.
\item Ils supplient le médecin de venir depuis ce matin.
\item Dès que tu auras fini, nous demanderons l'addition.
\item L'élève sollicite une bourse depuis deux mois, mais il a beau insister, on ne répond pas.
\end{enumerate}
\end{exercice}

\begin{corrige}[Thème]
\begin{enumerate}
\item \esp{Por más que trabaje, fracasa siempre.} (subjonctif après \esp{por más que})
\item \esp{Pido arroz en el mercado de Mokolo.} (\esp{pedir} + chose)
\item \esp{Le pregunta la hora al conductor del bendskin.} (\esp{preguntar} + information)
\item \esp{Vivimos en Yaoundé desde 2015.} (date : \esp{desde} + présent)
\item \esp{Mi padre vende cacao desde hace veinte años.} (durée : \esp{desde hace} ; autre solution : \esp{Hace veinte años que mi padre vende cacao.})
\item \esp{Desde que llegó el periodista, todo el mundo pide información.} (\esp{desde que} + passé : indicatif)
\item \esp{Por mucho que hable el ministro, la multitud no lo escucha.}
\item \esp{Ruegan al médico que venga desde esta mañana.} (\esp{rogar que} + subjonctif : \esp{venga})
\item \esp{Desde que hayas terminado, pediremos la cuenta.} (\esp{desde que} + futur : subjonctif ; solution plus courante : \esp{Cuando hayas terminado, pediremos la cuenta.})
\item \esp{El alumno solicita una beca desde hace dos meses, pero por más que insista, no responden.}
\end{enumerate}
\end{corrige}

\courssub{Entraînement à la version (ES $\rightarrow$ FR) : huit phrases justifiées}

\begin{exercice}[Version : traduisez en français et justifiez]
\begin{enumerate}
\item \esp{Por más que llame, nadie contesta en la redacción.}
\item \esp{El vendedor me preguntó por mi familia.}
\item \esp{Los alumnos piden silencio desde hace diez minutos.}
\item \esp{Desde que lleguemos a Douala, te llamaremos.}
\item \esp{La empresa solicita periodistas con experiencia.}
\item \esp{Te ruego que me esperes en la entrada.}
\item \esp{Aunque tiene mucho dinero, no es feliz.}
\item \esp{Hace tres años que mi hermana estudia español.}
\end{enumerate}
\end{exercice}

\begin{corrige}[Version]
\begin{enumerate}
\item « Il a beau appeler, personne ne répond à la rédaction. » — \esp{por más que} + subjonctif = « avoir beau » ; \esp{redacción} = « rédaction » (salle des journalistes), à ne pas confondre avec la « rédaction » scolaire, qui se dit \esp{redacción} aussi, mais le contexte décide.
\item « Le vendeur m'a demandé des nouvelles de ma famille. » — \esp{preguntar por} = demander des nouvelles de ; passé simple $\rightarrow$ passé composé.
\item « Les élèves demandent le silence depuis dix minutes. » — \esp{desde hace} + durée = « depuis » + durée ; présent dans les deux langues.
\item « Dès que nous arriverons à Douala, nous t'appellerons. » — \esp{desde que} + subjonctif = futur en français.
\item « L'entreprise recherche des journalistes expérimentés. » — \esp{solicitar} : demande formelle, ici « recherche, recrute » ; registre soutenu.
\item « Je te prie de m'attendre à l'entrée. » — \esp{rogar que} + subjonctif = prier, supplier ; \esp{esperes} = subjonctif présent.
\item « Même s'il a beaucoup d'argent, il n'est pas heureux. » — \esp{aunque} + indicatif signale un fait réel, constaté ; en français, « bien que » exigeant le subjonctif, la tournure « même si » + indicatif rend mieux la nuance de réalité.
\item « Ma sœur étudie l'espagnol depuis trois ans. » — \esp{hace\ldots que} + présent = « depuis » + durée ; ne pas traduire par un passé composé (« il y a trois ans qu'elle a étudié » serait fautif ici).
\end{enumerate}
\end{corrige}

\courssub{Auto-évaluation et préparation à l'activité d'intégration}

\begin{propriete}[Grille d'auto-évaluation de la traduction]
Après chaque phrase, cochez :
\begin{itemize}
\item[$\square$] \textbf{Sens exact} : le message français est fidèle, rien n'est ajouté ni omis.
\item[$\square$] \textbf{Verbe / structure correct} : \esp{pedir}/\esp{preguntar} bien choisi ; \esp{desde}/\esp{desde hace} juste.
\item[$\square$] \textbf{Mode correct} : subjonctif après \esp{por más que}, \esp{rogar que}, \esp{desde que} + futur.
\item[$\square$] \textbf{Orthographe et accents} : \esp{información}, \esp{está}, \esp{teléfono}, \esp{explicación}.
\item[$\square$] \textbf{Registre adapté} : \esp{solicitar} et \esp{rogar} réservés au soutenu ; \esp{pedir}/\esp{preguntar} au courant.
\end{itemize}
Cinq cases cochées : traduction excellente. Trois ou moins : reprendre la règle correspondante.
\end{propriete}

\begin{savaistu}[Traduction au baccalauréat]
Au Bac A, les items de traduction portent souvent sur des phrases courtes et isolées, exactement comme celles de cette leçon. En version, l'examinateur attend une reformulation naturelle en bon français, pas un calque mot à mot : \esp{desde hace una hora} doit devenir « depuis une heure », et non « depuis fait une heure ». Une phrase simple bien traduite vaut autant qu'une phrase difficile maladroitement calquée.
\end{savaistu}

\begin{experience}[Préparation : dialogue journaliste / vendeur]
Par binômes, préparez un dialogue entre un journaliste de passage dans un quartier de Yaoundé et un vendeur du marché. Contraintes obligatoires :
\begin{itemize}
\item une phrase avec \esp{por más que} ou \esp{por mucho que} + subjonctif ;
\item deux emplois de \esp{pedir} et deux de \esp{preguntar} (dont un \esp{preguntar por}) ;
\item une phrase avec \esp{desde} + date, une avec \esp{desde hace} + durée, une avec \esp{desde que} ;
\item un emploi de \esp{rogar que} + subjonctif pour une supplication polie.
\end{itemize}
Amorce possible : \esp{— Buenos días, señor. Soy periodista. ¿Puedo preguntarle por la vida del barrio? — Pídame lo que quiera, pero por más que pregunte, no todos le responderán\ldots}
\end{experience}

Cette synthèse achève notre étude des trois notions. Gardez le tableau récapitulatif à portée de main : c'est votre boîte à outils pour toute épreuve de thème et de version.

\newpage

\coursec{Activité d'intégration}

\coursec{Leçon E25 : Traduction : « avoir beau » ; « demander » ; « depuis »}

\courssub{Observation guidée des trois notions}

\begin{textebox}[Dialogue au marché Mokolo — Observation]
— Buenos días, don Félix. ¿\textbf{Desde cuándo} tiene usted este puesto de teléfonos?\par
— \textbf{Desde hace} quince años, \textbf{desde que} terminé mis estudios.\par
— ¿Y \textbf{qué les pide} a sus clientes?\par
— Les \textbf{pido} que paguen a tiempo. A veces les \textbf{pregunto} si quieren garantía.\par
— ¿Ha tenido que \textbf{rogar} a alguien?\par
— Sí, \textbf{por más que} le \textbf{rogué}, no me pagó. Tuve que \textbf{solicitar} la ayuda de un mediador.\par
— \textbf{Por mucho que} sea difícil, usted sigue aquí.\par
— Claro, \textbf{hace} quince años \textbf{que} trabajo aquí y no me voy.
\end{textebox}

\begin{definition}[Glossaire du dialogue]
\esp{don Félix} : Monsieur Félix (titre de respect) ; \esp{el puesto} : l'étal, le stand ; \esp{la garantía} : la garantie ; \esp{el mediador} : le médiateur ; \esp{no me voy} : je ne pars pas (verbe \esp{irse}).
\end{definition}

\courssub{Traduire « avoir beau » : concession et insistance}

\begin{propriete}[« Avoir beau » = Concession insistant sur l'inutilité de l'effort]
\emph{Structure de base :} \textbf{\esp{Por más que} / \esp{Por mucho que} + SUBJONCTIF} (le fait est réel, mais le résultat n'en découle pas).\par
\emph{Variante avec \esp{aunque} :}
\begin{itemize}
\item \esp{Aunque} + \textbf{SUBJONCTIF} : le fait est hypothétique, non réalisé ou mis en doute (\emph{« même si »}). Ex: \esp{Aunque llueva, iré.} « Même s'il pleut, j'irai. »
\item \esp{Aunque} + \textbf{INDICATIF} : le fait est réel, connu, admis (\emph{« bien que »}). Ex: \esp{Aunque llueve, salgo.} « Bien qu'il pleuve (il pleut vraiment), je sors. »
\end{itemize}
\emph{Autres formes :} \esp{Por} + adjectif/adverbe + \esp{que} + SUBJONCTIF (\esp{Por rápido que corras, no llegarás}).
\end{propriete}

\begin{exemplebox}[Exemples traduits — « Avoir beau »]
\esp{Por más que estudia, no aprende.} « Il a beau étudier, il n'apprend pas. » (Présent)\par
\esp{Por mucho que trabajé, no gané suficiente.} « J'ai beau avoir travaillé, je n'ai pas gagné assez. » (Passé : subj. imparfait \esp{trabajara/trabajase})\par
\esp{Por más que le ruegues, no vendrá.} « Tu as beau le supplier, il ne viendra pas. » (Subjonctif présent)\par
\esp{Aunque es rico, vive modestamente.} « Bien qu'il soit riche (fait réel), il vit modestement. » (Indicatif)\par
\esp{Aunque sea rico, no es feliz.} « Quand bien même il serait riche (hypothétique), il ne serait pas heureux. » (Subjonctif)
\end{exemplebox}

\begin{attention}[Piège majeur : le subjonctif est OBLIGATOIRE après \esp{por más que / por mucho que}]
\begin{itemize}
\item \textbf{FAUX :} \esp{*Por más que trabajo, no gano.} (Indicatif)
\item \textbf{VRAI :} \esp{Por más que \textbf{trabaje}, no gano.} (Subjonctif présent)
\item \textbf{FAUX :} \esp{*Por mucho que preparé la solicitud.} (Indicatif passé)
\item \textbf{VRAI :} \esp{Por mucho que \textbf{preparara/preparase} la solicitud.} (Subjonctif imparfait)
\end{itemize}
\emph{Rappel :} Le subjonctif marque ici la \textbf{déconnexion} entre l'effort (réel) et le résultat (négatif).
\end{attention}

\courssub{Traduire « demander » : cuatro verbos, cuatro matices}

\begin{propriete}[Les quatre verbes de la demande : construction et registre]
\begin{center}
\begin{tabularx}{\linewidth}{|X|X|X|X|}
\hline
\rowcolor{popblueL} Verbe & Sens précis & Construction type & Registre / Contexte \\
\hline
\textbf{\esp{PEDIR}} & Demander une \textbf{chose}, un service, un objet & \esp{pedir algo (a alguien)} & Courant, concret \\
\hline
\textbf{\esp{PREGUNTAR}} & Demander une \textbf{information}, poser une question & \esp{preguntar algo / por algo / si...} & Neutre, informationnel \\
\hline
\textbf{\esp{ROGAR}} & Supplier, prier, insister avec humilité & \esp{rogar a alguien que + SUBJ.} & Émotionnel, insistant \\
\hline
\textbf{\esp{SOLICITAR}} & Faire une demande \textbf{officielle, administrative, écrite} & \esp{solicitar algo (formal)} & Administratif, formel \\
\hline
\end{tabularx}
\end{center}
\end{propriete}

\begin{propriete}[Conjugaisons irrégulières à connaître (Présent de l'indicatif et Subjonctif présent)]
\begin{center}
\begin{tabular}{|l|c|c|c|c|}
\hline
\rowcolor{poporangeL} & \textbf{yo} & \textbf{tú/él} & \textbf{nosotros} & \textbf{ellos} \\
\hline
\textbf{PEDIR} (e$\to$i) & pido, pida & pides, pida & pedimos, pidamos & piden, pidan \\
\hline
\textbf{ROGAR} (o$\to$ue) & ruego, ruegue & ruegas, ruegue & rogamos, roguemos & ruegan, rueguen \\
\hline
\textbf{PREGUNTAR} (régulier) & pregunto, pregunte & preguntas, pregunte & preguntamos, preguntemos & preguntan, pregunten \\
\hline
\textbf{SOLICITAR} (régulier) & solicito, solicite & solicitas, solicite & solicitamos, solicitemos & solicitan, soliciten \\
\hline
\end{tabular}
\end{center}
\emph{Note :} Au passé, \esp{pedir} et \esp{rogar} changent de radical au Prétérit indéfini (3e pers.) : \esp{pidió, pidieron} ; \esp{rogó, rogaron}. Au Subjonctif imparfait : \esp{pidiera/pidiese}, \esp{rogara/rogase}.
\end{propriete}

\begin{exemplebox}[Exemples contrastés — Ne pas confondre !]
\esp{Le pido un favor.} « Je lui demande un service. » (\textbf{chose})\par
\esp{Le pregunto la hora.} « Je lui demande l'heure. » (\textbf{information})\par
\esp{Le ruego que me escuche.} « Je le supplie de m'écouter. » (\textbf{insistance + subjonctif})\par
\esp{He solicitado un visado.} « J'ai demandé un visa (dossier officiel). » (\textbf{administratif})\par
\esp{¿Me puedes pedir perdón?} « Peux-tu me demander pardon ? » (\textbf{chose abstraite = pedir})\par
\esp{¿Me puedes preguntar por qué?} « Peux-tu me demander pourquoi ? » (\textbf{information = preguntar})
\end{exemplebox}

\begin{attention}[Erreurs fréquentes des francophones]
\begin{itemize}
\item \textbf{Confusion \esp{pedir} / \esp{preguntar} :} On ne dit jamais \esp{*pregunto un favor}. \esp{Preguntar} ne prend \textbf{jamais} de COD « chose » ni de \esp{a} + personne pour l'objet demandé.
\item \textbf{\esp{Rogar} + Subjonctif :} \esp{Te ruego que \textbf{ven\textbf{g}as}} (et non \esp{*vienes}).
\item \textbf{\esp{Solicitar} vs \esp{Pedir} :} Au marché, on \esp{pide} un prix ; à la mairie, on \esp{solicita} une licence. \esp{Solicitar} implique une procédure.
\item \textbf{Faux ami :} \esp{Demandar} = « poursuivre en justice, intenter un procès » (terme juridique), \textbf{pas} « demander » au sens courant.
\end{itemize}
\end{attention}

\courssub{Traduire « depuis » : repères temporels et aspect}

\begin{propriete}[Les quatre structures de l'antériorité / durée]
\begin{center}
\begin{tabularx}{\linewidth}{|X|X|X|}
\hline
\rowcolor{popgreenL} Structure & Valeur & Exemple \\
\hline
\textbf{\esp{DESDE} + nom/date/heure} & Point de départ \textbf{ponctuel} dans le passé & \esp{Desde 2010}, \esp{desde el lunes}, \esp{desde mi llegada} \\
\hline
\textbf{\esp{DESDE HACE} + durée} & \textbf{Durée} écoulée depuis le début (action encore en cours) & \esp{Desde hace dos años}, \esp{desde hace mucho tiempo} \\
\hline
\textbf{\esp{HACE} + durée + \esp{QUE} + verbe (PRÉSENT)} & \textbf{Durée} + action \textbf{en cours} (insiste sur la continuité) & \esp{Hace dos años que vivo aquí.} \\
\hline
\textbf{\esp{DESDE QUE} + verbe (INDICATIF)} & Point de départ \textbf{exprimé par une action passée} (antériorité) & \esp{Desde que llegué}, \esp{desde que era niño} \\
\hline
\end{tabularx}
\end{center}
\end{propriete}

\begin{methode}[Choisir la bonne structure : arbre de décision]
\begin{enumerate}
\item La référence temporelle est-elle une \textbf{date / un moment précis} ? $\to$ \esp{Desde 2020}, \esp{desde ayer}.
\item Est-ce une \textbf{durée} (\emph{depuis 2 ans}) ?
    \begin{itemize}
    \item Si focus sur la \textbf{durée} elle-même $\to$ \esp{Desde hace 2 años} / \esp{Hace 2 años que...}
    \item Si focus sur le \textbf{début de l'action} $\to$ \esp{Desde hace 2 años} (synonyme).
    \end{itemize}
\item Le point de départ est-il une \textbf{action passée} (\emph{depuis que je suis arrivé}) ? $\to$ \esp{Desde que llegué} (Indicatif passé).
\item \textbf{Aspect verbal :}
    \begin{itemize}
    \item Action \textbf{toujours en cours} $\to$ \textbf{Présent} : \esp{Desde hace 2 años \textbf{estudio} / \textbf{estudio desde hace 2 años}.}
    \item Action \textbf{terminée} dans le passé $\to$ \textbf{Prétérit} : \esp{Estudié desde 2010 hasta 2015.}
    \item Antériorité dans un récit au passé $\to$ \textbf{Plus-que-parfait} : \esp{Llegué a las 8. Desde las 6 \textbf{había estado} esperando.}
    \end{itemize}
\end{enumerate}
\end{methode}

\begin{exemplebox}[Variantes stylistiques — « Depuis »]
\esp{Vivo en Yaoundé desde 2015.} (Point de départ)\par
\esp{Vivo en Yaoundé desde hace nueve años.} (Durée)\par
\esp{Hace nueve años que vivo en Yaoundé.} (Durée + Présent = très courant à l'oral)\par
\esp{Vivo en Yaoundé desde que me casé.} (Action passée comme repère)\par
\esp{\textbf{¡Ojo!} \esp{Nunca} \esp{*desde nueve años}. Toujours \esp{desde hace} ou \esp{hace... que}.}
\end{exemplebox}

\begin{attention}[Pièges « Depuis » — Spécifiques francophones]
\begin{itemize}
\item \textbf{*Desde dos años} $\to$ \textbf{INCORRECT}. Le français « depuis + durée » trompe. Espagnol : \esp{DESDE HACE} + durée OU \esp{HACE} + durée + \esp{QUE}.
\item \textbf{Depuis que + Subjonctif ?} $\to$ \esp{Desde que} + \textbf{INDICATIF} (passé ou présent) pour un fait réel accompli. Subjonctif seulement si l'action est \textbf{future/anticipée} : \esp{Avísame desde que \textbf{llegue}} (Quand tu auras arrivé).
\item \textbf{\esp{Desde} vs \esp{De} :} \esp{De lunes a viernes} (De lundi à vendredi = intervalle), \esp{Desde el lunes} (Depuis lundi = point de départ).
\item \textbf{Accent sur \esp{está} / \esp{esta} / \esp{hace} :} \esp{Está aquí desde hace una hora} (il est), \esp{Esta casa} (cette maison), \esp{Hace una hora} (il y a une heure).
\end{itemize}
\end{attention}

\courssub{Synthèse comparative, pièges et entraînement thème/version}

\courssub{Objectifs de l'activité}

À l'issue de cette activité d'intégration, l'élève sera capable de :
\begin{itemize}
\item mobiliser dans une production orale et écrite authentique les trois notions étudiées : la traduction de « avoir beau » (\esp{por más que}, \esp{por mucho que}, \esp{aunque} + subjonctif), de « demander » (\esp{pedir}, \esp{preguntar}, \esp{rogar}, \esp{solicitar}) et de « depuis » (\esp{desde}, \esp{desde hace}, \esp{hace... que}, \esp{desde que}) ;
\item distinguer spontanément \esp{pedir} (demander une chose) et \esp{preguntar} (poser une question) dans un échange réel ;
\item employer correctement \esp{desde hace} + durée et \esp{desde que} + verbe pour exprimer l'ancienneté ;
\item tenir un rôle de communication (journaliste / commerçant) avec un registre de langue adapté ;
\item s'auto-corriger et corriger ses pairs à partir d'une grille ciblée sur les trois points de grammaire.
\end{itemize}

\courssub{Situation}

Tu es journaliste stagiaire à \emph{Radio Mvolyé}, une radio de proximité de Yaoundé. Dans le cadre de l'émission hebdomadaire \esp{« Voces del barrio »}, consacrée à la vie économique des quartiers, ta rédactrice en chef t'envoie au marché Mokolo réaliser un micro-trottoir en espagnol, destiné à être diffusé lors de la semaine hispanophone du lycée. Ta mission : interviewer un commerçant sur son parcours, ses relations avec les clients et les fournisseurs, et les difficultés qu'il rencontre malgré ses efforts.

Voici la fiche de mission remise par la rédactrice en chef :

\begin{textebox}[Ficha de misión — Radio Mvolyé]
\textbf{Emisión :} « Voces del barrio » — Especial mercado de Mokolo\par
\textbf{Tema del reportaje :} La vida de los comerciantes del mercado.\par
\textbf{Pistas para el periodista :}
\begin{itemize}
\item ¿Desde cuándo trabaja aquí el comerciante? ¿Cuánto tiempo hace que vende sus productos?
\item ¿Qué pide a sus clientes? ¿Qué les pregunta? ¿Ha rogado alguna vez a un cliente difícil?
\item ¿Ha presentado una solicitud oficial (ayuda del ayuntamiento, crédito bancario)?
\item ¿Le pasa que, por más que se esfuerza, los resultados no llegan?
\end{itemize}
\textbf{Duración :} 2 minutos de audio o diálogo escrito de 150 palabras.\par
\textbf{Lema de la emisión :} « Por más que la vida es dura, el mercado nunca duerme. »
\end{textebox}

\begin{definition}[Glossaire de la fiche de mission]
\esp{la emisión} : l'émission ; \esp{las pistas} : les pistes, les indications ; \esp{pedir} : demander (une chose) ; \esp{preguntar} : demander (une information) ; \esp{rogar} : supplier ; \esp{la solicitud} : la demande officielle, le dossier ; \esp{por más que} : avoir beau ; \esp{esforzarse} : faire des efforts ; \esp{el ayuntamiento} : la mairie ; \esp{el lema} : la devise.
\end{definition}

\courssub{Consignes}

\textbf{Tâche 1 — Compréhension et repérage (10 min).} Lis la fiche de mission. Souligne en bleu les trois verbes qui expriment une demande, en vert les deux expressions de l'ancienneté, en rouge l'expression de la concession. Recopie-les dans ton cahier en complétant ce tableau :

\begin{center}
\begin{tabularx}{\linewidth}{|X|X|X|}
\hline
\rowcolor{popblueL} Notion française & Structure espagnole & Exemple de la fiche \\
\hline
« demander » (une chose) & \esp{pedir} & \esp{¿Qué pide a sus clientes?} \\
\hline
« demander » (une information) & \esp{preguntar} & \esp{¿Qué les pregunta?} \\
\hline
« supplier » & \esp{rogar} & \esp{¿Ha rogado alguna vez...?} \\
\hline
« depuis » (+ durée) & \esp{hace... que} / \esp{desde hace} & \esp{¿Cuánto tiempo hace que vende...?} \\
\hline
« avoir beau » & \esp{por más que} + subj. & \esp{por más que se esfuerza} \\
\hline
\end{tabularx}
\end{center}

\textbf{Tâche 2 — Préparation du journaliste (10 min).} En binôme, l'élève journaliste rédige \textbf{cinq questions} en espagnol, chacune mobilisant une structure imposée :
\begin{enumerate}
\item une question sur l'ancienneté avec \esp{¿Desde cuándo...?} ou \esp{¿Cuánto tiempo hace que...?} ;
\item une question avec \esp{pedir} (ce que le commerçant demande aux clients ou aux fournisseurs) ;
\item une question avec \esp{preguntar} (une information qu'il demande) ;
\item une question avec \esp{rogar} ou \esp{solicitar} (un client difficile, une démarche officielle) ;
\item une question avec \esp{por más que} + subjonctif (des efforts vains).
\end{enumerate}

\textbf{Tâche 3 — Jeu de rôle (10 min).} L'élève commerçant choisit un étal (mangues, tissus pagnes, téléphones et accessoires, légumes) et répond \textbf{spontanément}, en veillant à réutiliser au moins une fois chaque structure : \esp{desde hace} + durée, \esp{desde que} + verbe, \esp{pedir}, \esp{preguntar}, \esp{rogar} ou \esp{solicitar}, et une phrase de concession avec \esp{por mucho que} ou \esp{aunque} + subjonctif.

\textbf{Tâche 4 — Production (15 min).} Deux options au choix :
\begin{itemize}
\item \textbf{Option orale :} enregistre l'interview (2 minutes maximum) avec le téléphone du groupe ; introduis-la par une phrase de présentation (\esp{« Hablamos con... desde el mercado de Mokolo »}).
\item \textbf{Option écrite :} rédige le dialogue complet (environ 150 mots), avec incise de présentation et conclusion du journaliste.
\end{itemize}

\textbf{Tâche 5 — Partage et correction collective (15 min).} Deux binômes présentent leur production. La classe complète la grille d'écoute ci-dessous, puis le professeur corrige collectivement en ciblant uniquement les trois notions de la leçon.

\begin{center}
\begin{tabularx}{\linewidth}{|X|X|X|}
\hline
\rowcolor{popgreenL} Critère visé & Réussi (oui/non) & Phrase entendue \\
\hline
\esp{desde hace} + durée correctement employé & & \\
\hline
\esp{desde que} + verbe conjugué (indicatif) & & \\
\hline
\esp{pedir} et \esp{preguntar} non confondus & & \\
\hline
\esp{por más/mucho que} + subjonctif & & \\
\hline
\end{tabularx}
\end{center}

\courssub{Ressources à mobiliser}

\begin{aretenir}[Boîte à outils de la leçon — Récapitulatif synthétique]
\begin{itemize}
\item \textbf{« Avoir beau » :} \esp{por más que} / \esp{por mucho que} + \textbf{subjonctif} : \esp{Por más que \textbf{grite}, nadie me escucha.} « J'ai beau crier, personne ne m'écoute. » Avec \esp{aunque} : subjonctif si hypothétique, indicatif si réel.
\item \textbf{« Demander » :} \esp{pedir algo a alguien} (chose), \esp{preguntar algo / si...} (information), \esp{rogar a alguien que + subj.} (supplier), \esp{solicitar algo} (officiel : \esp{solicitar un crédito}).
\item \textbf{« Depuis » :} \esp{desde} + date (\esp{desde 2015}) ; \esp{desde hace} + durée (\esp{desde hace diez años}) ; \esp{desde que} + verbe indicatif (\esp{desde que llegué}) ; \esp{hace diez años que vendo mangas} (présent = en cours).
\end{itemize}
\end{aretenir}

\begin{attention}[Pièges à éviter pendant le jeu de rôle — Rappel ciblé]
\begin{itemize}
\item Ne dis jamais \esp{*desde diez años} pour « depuis dix ans » : avec une durée, il faut \esp{desde hace diez años} ou \esp{hace diez años que...}.
\item \esp{Preguntar} ne se construit jamais avec une personne introduite par \esp{a} pour la chose demandée : on dit \esp{pedir un favor a un cliente}, mais \esp{preguntar el precio}.
\item Après \esp{por más que}, le subjonctif est \textbf{obligatoire} : \esp{por más que \textbf{trabaje}} (et non \esp{*trabajo}).
\item Attention à l'accent : \esp{está} (il est, verbe \esp{estar}) / \esp{esta} (cette, démonstratif) ; et à l'orthographe de \esp{hace} (temps, verbe \esp{hacer}) sans confusion avec \esp{haze} (inexistant).
\item \esp{Rogar} (o$\to$ue) : \esp{yo ruego, tú ruegas, él ruega}... Subjonctif : \esp{que yo ruegue}.
\end{itemize}
\end{attention}

\courssub{Proposition de production et corrigé commenté}

\begin{exoresolu}[Modèle d'interview : « Voces del mercado »]
Rédige le dialogue complet de l'interview (environ 150 mots) en mobilisant les trois notions de la leçon : « avoir beau », « demander » (quatre verbes) et « depuis ».
\tcblower
\textbf{Production modèle :}

\begin{textebox}[Voces del mercado — entrevista con doña Solange]
— Buenos días, oyentes de Radio Mvolyé. Hablamos hoy con doña Solange, vendedora de mangos, desde el mercado de Mokolo. Doña Solange, ¿desde cuándo trabaja usted aquí?\par
— Trabajo en este puesto \textbf{desde hace} veinte años, \textbf{desde que} llegué de mi pueblo.\par
— ¡Veinte años! ¿Qué les \textbf{pide} usted a sus clientes?\par
— Solo les \textbf{pido} un poco de paciencia cuando hay mucha gente. Y les \textbf{pregunto} siempre si quieren mangos maduros o verdes.\par
— ¿Ha \textbf{rogado} alguna vez a un cliente difícil?\par
— Sí, una vez le \textbf{rogué} a un cliente que me pagara, pero \textbf{por más que} le \textbf{rogara}, no quiso pagarme.\par
— ¿Ha \textbf{solicitado} ayuda al ayuntamiento?\par
— Sí, \textbf{solicité} un crédito hace dos años, pero \textbf{por mucho que} \textbf{preparara} mi solicitud, no me la aceptaron.\par
— ¿Y la publicidad?\par
— \textbf{Por más que} \textbf{haga} publicidad, los clientes nuevos no vienen. Pero no me rindo: \textbf{hace} veinte años \textbf{que} vendo mangos y seguiré aquí.\par
— Gracias, doña Solange. Desde Mokolo, para « Voces del barrio », les habló Radio Mvolyé.
\end{textebox}

\textbf{Commentaire des choix de langue :}
\begin{enumerate}
\item \esp{¿Desde cuándo trabaja usted aquí?} : question d'ancienneté classique ; le vouvoiement (\esp{usted}) convient au registre journalistique.
\item \esp{desde hace veinte años} : durée, donc \esp{desde hace} obligatoire ; \esp{desde que llegué} : point de départ exprimé par une proposition verbale au passé (indicatif).
\item \esp{les pido paciencia} / \esp{les pregunto si...} : la distinction est respectée — \esp{pedir} pour la chose demandée (la patience), \esp{preguntar} pour l'information (mûres ou vertes).
\item \esp{le rogué... que me pagara} : \esp{rogar} + subordonnée au subjonctif imparfait, concordance des temps correcte après un passé simple (\esp{rogué}).
\item \esp{por más que le \textbf{rogara}, no quiso pagarme} : concession au passé avec \textbf{subjonctif imparfait} ; \esp{por mucho que \textbf{preparara} mi solicitud} : même structure, variante lexicale (\esp{por mucho que}).
\item \esp{solicité un crédito} / \esp{mi solicitud} : verbe et nom de la famille de la demande officielle, bien distingués de \esp{pedir}.
\item \esp{hace veinte años que vendo} : troisième tournure possible de « depuis », avec présent de valeur durative — excellente variante stylistique pour conclure.
\item \esp{Por más que \textbf{haga} publicidad} : concession au présent, subjonctif présent (\esp{haga} de \esp{hacer}).
\end{enumerate}
\end{exoresolu}

\courssub{Bilan de l'activité}

Cette activité a permis de réinvestir les trois notions de la leçon dans une situation de communication authentique et ancrée dans le quotidien camerounais. Pour vérifier que les objectifs sont atteints, chaque élève coche la grille d'auto-évaluation :

\begin{center}
\begin{tabularx}{\linewidth}{|X|X|}
\hline
\rowcolor{popgoldL} Je sais... & Oui / Encore fragile \\
\hline
traduire « avoir beau » par \esp{por más que} / \esp{por mucho que} + subjonctif & \\
\hline
choisir entre \esp{pedir}, \esp{preguntar}, \esp{rogar} et \esp{solicitar} selon le sens & \\
\hline
employer \esp{desde}, \esp{desde hace}, \esp{desde que} et \esp{hace... que} sans confusion & \\
\hline
tenir un dialogue de 150 mots avec un registre adapté & \\
\hline
\end{tabularx}
\end{center}

\begin{methode}[Pour aller plus loin : réutiliser la grille en version]
\begin{enumerate}
\item Reprends ton dialogue et traduis-le en français : c'est un excellent exercice de version ciblée.
\item Inverse les rôles : le commerçant devient journaliste et interviewe un client sur ses habitudes d'achat (\esp{¿Desde cuándo compra aquí? ¿Qué pide siempre?}).
\item À la maison, rédige un court article de presse (100 mots) à partir de l'interview, en passant du style direct au style indirect : \esp{Doña Solange dijo que trabajaba allí desde hacía veinte años...} (Notez le passage \esp{desde hace} $\to$ \esp{desde hacía} au discours indirect au passé).
\end{enumerate}
\end{methode}

\begin{savaistu}[Le marché, une école de langues]
À Mokolo comme sur les marchés de Malabo ou de Bata en Guinée équatoriale — seul pays hispanophone d'Afrique subsaharienne —, les commerçants jonglent chaque jour entre plusieurs langues : français, espagnol, langues locales. Savoir demander avec le bon verbe (\esp{pedir} un prix, \esp{preguntar} une information, \esp{rogar} en négociant) est une compétence de vie autant qu'une compétence grammaticale : au marché, celui qui choisit le bon mot obtient souvent le meilleur prix.
\end{savaistu}

\newpage

\coursec{Méthodes et exercices résolus}

\coursec{Leçon E25 : Traduction -- « avoir beau » ; « demander » ; « depuis »}

\courssub{Observation guidée : trois notions dans un contexte camerounais}

\begin{textebox}[Article de presse simulé -- Contexte : Yaoundé, médias et vie locale]
Desde hace tres meses, los habitantes del barrio Mvog-Ada \textbf{piden} la reparación del puente sobre el Mfoundi. Por mucho que \textbf{pregunten} a las autoridades, la respuesta siempre es la misma: « no hay presupuesto ». \textbf{Desde que} se derrumbó una parte del tablero, los taxistas \textbf{rogan} a los pasajeros que tengan cuidado. Una ONG local ha \textbf{solicitado} oficialmente ayuda al Ministerio de Obras Públicas. Los periodistas de la radio comunitaria \textbf{tienen beau} -- \emph{por más que} -- investigar, el expediente sigue cerrado.
\end{textebox}

\begin{experience}[Analyse linguistique du texte]
\begin{enumerate}
\item Relevez les cinq segments en gras et identifiez la notion française qu'ils traduisent (« demander », « depuis », « avoir beau »).
\item Pour \esp{pregunten} et \esp{rogan} : quel mode ? Pourquoi ?
\item Pour \esp{Desde hace tres meses} et \esp{Desde que se derrumbó} : quelle différence de structure ?
\item Traduisez le paragraphe en français naturel.
\end{enumerate}
\end{experience}

\begin{corrige}[Exercice d'observation]
\begin{enumerate}
\item \esp{piden} = demandent (une chose) ; \esp{pregunten} = posent des questions / demandent (une info) ; \esp{Desde hace tres meses} = cela fait trois mois que / depuis trois mois ; \esp{Desde que} = depuis que ; \esp{rogan} = supplient ; \esp{solicitado} = sollicité / demandé officiellement ; \esp{por más que investigan} (corrigé en \esp{investiguen} dans le texte final) = ont beau enquêter / quel que soit leur effort pour enquêter.
\item \esp{pregunten}, \esp{rogan} (subjonctif présent) : subordonnées après verbes de volonté/demande (\esp{preguntar que?} non, mais structure \esp{rogar que} + subj. ; ici \esp{pregunten} est subjonctif car dans une relative à antécédent indéfini/négatif implicite ? Non, ici \esp{pregunten} est à l'indicatif dans le texte original proposé par le collègue ? \textbf{Correction cruciale} : Dans le texte ci-dessus, \esp{pregunten} doit être à l'indicatif \esp{preguntan} car c'est un fait réel : « ils posent des questions ». \esp{rogan} est indicatif (fait réel). Le subjonctif n'apparaît que dans \esp{rogan \textbf{que} tengan cuidado}. \textbf{Le texte support a été corrigé ci-dessus pour refléter la réalité grammaticale : \esp{preguntan}, \esp{rogan} (indicatif), \esp{tengan} (subj).})
\item \esp{Desde hace} + durée + présent ; \esp{Desde que} + proposition conjuguée (ici passé simple \esp{derrumbó} car événement ponctuel passé).
\item « Cela fait trois mois que les habitants du quartier Mvog-Ada demandent la réparation du pont sur le Mfoundi. Quel que soit le nombre de fois où ils interrogent les autorités, la réponse est toujours la même : "pas de budget". Depuis qu'une partie du tablier s'est effondrée, les taximen supplient les passagers de faire attention. Une ONG locale a sollicité officiellement l'aide du Ministère des Travaux Publics. Les journalistes de la radio communautaire ont beau enquêter, le dossier reste classé. »
\end{enumerate}
\end{corrige}

\courssub{Traduire « avoir beau » : concession et insistance}

\begin{propriete}[Structure de la concession avec idée d'inutilité]
« Avoir beau + infinitif » exprime une concession réelle ou hypothétique teintée d'insistance ou d'inutilité. \textbf{Jamais} \esp{tener} + adjectif.
\begin{center}
\begin{tabularx}{\linewidth}{|X|X|X|}
\hline
\rowcolor{popblueL} \textbf{Structure espagnole} & \textbf{Mode} & \textbf{Exemple / Traduction} \\
\hline
\esp{Por mucho que} + verbe & Subjonctif (obligatoire) & \esp{Por mucho que llueva, saldremos.} (Quoiqu'il pleuve / Il a beau pleuvoir, nous sortirons.) \\
\hline
\esp{Por más que} + verbe & Subjonctif (obligatoire) & \esp{Por más que estudie, no aprueba.} (Il a beau étudier, il échoue.) \\
\hline
\esp{Por} + adj./adv. + \esp{que} + verbe & Subjonctif (obligatoire) & \esp{Por rico que sea, no es feliz.} (Il a beau être riche, il n'est pas heureux.) \\
\hline
\esp{Aunque} + verbe & Subjonctif (hypothétique/futur) / Indicatif (fait réel) & \esp{Aunque llueva, iremos.} (Même s'il pleut / Qu'il pleuve.) -- \esp{Aunque llueve, salgo.} (Comme il pleut / Bien qu'il pleuve -- fait constaté.) \\
\hline
\end{tabularx}
\end{center}
\end{propriete}

\begin{aretenir}[Concordance des temps avec « avoir beau »]
\begin{itemize}
\item Principale au \textbf{présent / futur} $\rightarrow$ Subordonnée au \textbf{subjonctif présent} : \esp{Por mucho que \textbf{intente}, no \textbf{consigue}.}
\item Principale au \textbf{passé (imparfait, passé composé, plus-que-parfait)} $\rightarrow$ Subordonnée au \textbf{subjonctif imparfait} (\esp{-ra} ou \esp{-se}) : \esp{Por más que \textbf{intentara/intentase}, no \textbf{consiguió}.}
\item \esp{Por} + adjectif + \esp{que} : l'adjectif s'accorde avec le sujet de la subordonnée. \esp{Por listo que \textbf{sea}...} (sujet \esp{él}) ; \esp{Por lista que \textbf{sea}...} (sujet \esp{ella}).
\end{itemize}
\end{aretenir}

\begin{methode}[Traduire « avoir beau » : démarche en 4 étapes]
\begin{enumerate}
\item \textbf{Identifier le sens} : concession + inutilité/insistance. Exclure \esp{tener}.
\item \textbf{Choisir la structure} selon la nature du mot-clé français :
\begin{itemize}
\item Verbe d'action $\rightarrow$ \esp{Por mucho que} / \esp{Por más que} + subjonctif.
\item Adjectif / Nom (être + qualité) $\rightarrow$ \esp{Por} + adj/nom + \esp{que} + subjonctif.
\item Registre plus neutre / « même si » $\rightarrow$ \esp{Aunque} + indicatif (fait) ou subjonctif (hypothèse).
\end{itemize}
\item \textbf{Accorder le mode et le temps} : vérifier le temps de la principale pour choisir subj. présent ou imparfait.
\item \textbf{Placer la subordonnée} (début ou fin de phrase) et vérifier les accents (\esp{está}, \esp{sí}, \esp{tú}...).
\end{enumerate}
\end{methode}

\begin{exoresolu}[Thème : « avoir beau » -- Niveau Bac]
Traduisez en espagnol :
1. Il a beau travailler jour et nuit, il ne gagne pas assez pour sa famille.
2. Elle a beau être riche, elle n'est pas heureuse.
3. Ils avaient beau crier, personne ne les entendait.
4. Quels que soient tes efforts, tu ne le convaincras pas. (Registre soutenu)
\tcblower
\textbf{Phrase 1.} Action (\esp{trabajar}), principale présent (\esp{no gana}) $\rightarrow$ Subj. présent. \esp{Por mucho que \textbf{trabaje} día y noche, no gana lo suficiente para su familia.} (\esp{lo suficiente} = assez).

\textbf{Phrase 2.} Attribut (\esp{rica}), principale présent (\esp{no es}) $\rightarrow$ \esp{Por} + adj + \esp{que} + subj. présent. \esp{Por \textbf{rica} que \textbf{sea}, no es feliz.} (Accord fém. \esp{rica}, subj. \esp{ser} $\rightarrow$ \esp{sea}).

\textbf{Phrase 3.} Action (\esp{gritar}), principale imparfait (\esp{oía}) $\rightarrow$ Subj. imparfait. \esp{Por más que \textbf{gritaran/gritasen}, nadie los \textbf{oía}.} (Pronom \esp{los} = les, personnes).

\textbf{Phrase 4.} « Quels que soient tes efforts » = concession sur nom/quantité. \esp{Por más \textbf{esfuerzos} que \textbf{hagas}, no lo \textbf{convencerás}.} (\esp{hagas} subj. présent \esp{hacer} ; futur \esp{convencerás} dans principale).
\end{exoresolu}

\courssub{Traduire « demander » : cuatro verbos, cuatro matices}

\begin{definition}[Les quatre piliers de la demande en espagnol]
\begin{description}
\item[\esp{Pedir}] Demander \textbf{une chose, un objet, un service, une permission, une action} (construction \esp{pedir algo a alguien} ou \esp{pedir que} + subjonctif). \emph{Verbe à diphtongue : \esp{pido, pides, pide, pedimos, pedís, piden} ; irrégulier au passé simple (\esp{pedí, pidió, pidieron}) et participe (\esp{pedido}).}
\item[\esp{Preguntar}] Demander \textbf{une information}, poser une question. \esp{Preguntar algo a alguien} / \esp{preguntar por algo} / \esp{preguntar si/qué/cuándo...} + indicatif (question directe/indirecte réelle).
\item[\esp{Rogar}] Supplier, prier, demander avec \textbf{insistance/émotion} (registre soutenu/littéraire). \esp{Rogar que} + \textbf{subjonctif}. \emph{Diphtongue : \esp{ruego, ruegas, ruega...} ; changement \esp{g $\rightarrow$ gu} : \esp{roguemos, roguéis}.}
\item[\esp{Solicitar}] Demander \textbf{officiellement, formellement} (emploi, bourse, visa, document, audience). Registre administratif/juridique. Transitif direct : \esp{solicitar algo}.
\end{description}
\end{definition}

\begin{attention}[Faux ami majeur : \esp{Demandar}]
\esp{Demandar} = \textbf{Poursuivre en justice, porter plainte}. \textbf{Jamais} « demander » au sens courant.
\esp{El vecino demandó al constructor por daños.} (Le voisin a poursuivi le constructeur pour dommages.)
\end{attention}

\begin{propriete}[Constructions types à mémoriser]
\begin{center}
\begin{tabularx}{\linewidth}{|l|X|X|}
\hline
\rowcolor{popblueL} \textbf{Verbe} & \textbf{Construction} & \textbf{Exemple} \\
\hline
\esp{Pedir} & \esp{pedir algo (a alguien)} & \esp{Le pedí un favor a mi hermano.} \\
\hline
\esp{Pedir} (action) & \esp{pedir que} + \textbf{subjonctif} & \esp{Te pido que \textbf{ven\textcolor{popred}{gas}}.} (Je te demande de venir.) \\
\hline
\esp{Preguntar} & \esp{preguntar algo / por algo / si...} & \esp{Preguntó \textbf{si} veníamos.} (Il a demandé si nous venions.) \\
\hline
\esp{Rogar} & \esp{rogar que} + \textbf{subjonctif} & \esp{Te \textbf{ruego} que me \textbf{escuches}.} \\
\hline
\esp{Solicitar} & \esp{solicitar algo (a alguien)} & \esp{Han \textbf{solicitado} una beca \textbf{al} ministerio.} \\
\hline
\end{tabularx}
\end{center}
\end{propriete}

\begin{exoresolu}[Version : « demander » -- Contexte médias/paix]
Traduisez en français : \esp{« El periodista \textbf{preguntó} al alcalde si el mercado de Yaundé estaría listo en diciembre. Los vecinos, por su parte, \textbf{pidieron} que se \textbf{arreglara} la carretera, y una ONG \textbf{solicitó} oficialmente ayuda al ayuntamiento. »}
\tcblower
\textbf{Analyse.}
\begin{itemize}
\item \esp{preguntó} (info) $\rightarrow$ « demanda / pose la question ».
\item \esp{pidieron que se arreglara} (action/réparation) $\rightarrow$ « demandèrent que » + subjonctif français (\esp{se arreglara} = subj. imparfait $\rightarrow$ \esp{soit réparée / fût réparée}).
\item \esp{solicitó} (officiel) $\rightarrow$ « sollicita / demande officiellement ».
\item \esp{ayuntamiento} = mairie / conseil municipal (faux ami : $\neq$ ajustement).
\item \esp{estaría} (conditionnel) = futur dans le passé $\rightarrow$ « serait ».
\end{itemize}
\textbf{Traduction.} « Le journaliste demanda au maire si le marché de Yaundé serait prêt en décembre. Les voisins, de leur côté, demandèrent que la route soit réparée, et une ONG sollicita officiellement l'aide de la mairie. »
\end{exoresolu}

\begin{exoresolu}[Thème : « demander » -- Pièges]
Traduisez :
1. Je te demande d'écouter. 2. Ils ont demandé des informations sur le train. 3. Nous sollicitons un rendez-vous avec le proviseur. 4. Il supplia son père de lui pardonner.
\tcblower
1. \esp{Te \textbf{pido} que \textbf{escuches}.} (Action $\rightarrow$ \esp{pedir que} + subj.). \textbf{Pas} \esp{*te pido escuchar}.
2. \esp{Pidieron información sobre el tren.} (Chose $\rightarrow$ \esp{pedir}). \textbf{Pas} \esp{*preguntaron información}.
3. \esp{\textbf{Solicitamos} una cita con el director.} (Officiel $\rightarrow$ \esp{solicitar}).
4. \esp{\textbf{Rogó} a su padre que le \textbf{perdonara}.} (Insistance + passé $\rightarrow$ \esp{rogar que} + subj. imparfait).
\end{exoresolu}

\courssub{Traduire « depuis » : repères temporels et aspect}

\begin{propriete}[Les trois structures du « depuis » temporel]
\begin{center}
\begin{tabularx}{\linewidth}{|X|X|X|}
\hline
\rowcolor{popblueL} \textbf{Français} & \textbf{Espagnol} & \textbf{Règle clé} \\
\hline
Depuis + \textbf{point de départ} (date, événement) & \esp{Desde} + indicatif (présent, passé simple, imparfait...) & \esp{Vivo en Duala \textbf{desde} 2015.} / \esp{\textbf{Desde} que \textbf{llegó}, todo cambió.} \\
\hline
Depuis + \textbf{durée} (action qui continue) & \textbf{A.} \esp{Desde hace} + durée + \textbf{présent} \newline \textbf{B.} \esp{Hace} + durée + \esp{que} + \textbf{présent} & \esp{Estudio español \textbf{desde hace} tres años.} \newline \esp{\textbf{Hace} tres años \textbf{que} estudio español.} \\
\hline
« Il y a ... que » / « Cela fait ... que » (passé, action terminée ou continue dans le passé) & \esp{Hacía} + durée + \esp{que} + \textbf{imparfait} & \esp{\textbf{Hacía} dos años \textbf{que} \textbf{vivía} allí.} (Il y habitait depuis deux ans / Cela faisait deux ans qu'il y habitait.) \\
\hline
\end{tabularx}
\end{center}
\end{propriete}

\begin{aretenir}[Le piège du passé composé -- Règle d'or]
\textbf{Jamais} de \esp{Pretérito Perfecto Compuesto} (\esp{he vivido}) avec \esp{desde/hace} si l'action \textbf{continue encore au présent}.
\begin{itemize}
\item \esp{*He vivido aquí desde hace cinco años.} $\leftarrow$ \textbf{FAUX} (calque français « j'habite ici depuis 5 ans »).
\item \esp{Vivo aquí desde hace cinco años.} / \esp{Hace cinco años que vivo aquí.} $\leftarrow$ \textbf{CORRECT} (Présent = action en cours).
\end{itemize}
Le \esp{Pretérito Perfecto} s'utilise si l'action est \textbf{terminée} : \esp{He estado enfermo desde el lunes hasta el viernes.} (J'ai été malade du lundi au vendredi).
\end{aretenir}

\begin{methode}[Choisir la bonne structure « depuis »]
\begin{enumerate}
\item \textbf{Analyser le complément} : Date/Événement $\rightarrow$ \esp{Desde}. Durée $\rightarrow$ \esp{Desde hace} / \esp{Hace... que}. Proposition (\esp{depuis que}) $\rightarrow$ \esp{Desde que}.
\item \textbf{Déterminer l'aspect} : Action \textbf{toujours en cours} $\rightarrow$ \textbf{Présent} espagnol. Action \textbf{terminée dans le passé} $\rightarrow$ \textbf{Imparfait} espagnol (\esp{Hacía... que} + imparfait).
\item \textbf{Vérifier le mode après \esp{Desde que}} : \textbf{Indicatif} (fait réel, certain). \esp{Desde que \textbf{sé} la verdad...} / \esp{Desde que \textbf{llegó}...}. Pas de subjonctif.
\end{enumerate}
\end{methode}

\begin{exoresolu}[Thème : « depuis » -- Variété structurelle]
Traduisez :
1. Nous apprenons l'espagnol depuis la classe de seconde.
2. Il pleut depuis trois jours à Yaoundé.
3. Depuis qu'elle travaille à la radio, elle est célèbre.
4. Il y avait longtemps qu'ils attendaient le bus.
\tcblower
\textbf{1. Point de départ.} \esp{Aprendemos español \textbf{desde} la clase de segundo.} (ou \esp{desde segundo de secundaria}).
\textbf{2. Durée + présent.} Deux solutions équivalentes : \esp{Llueve \textbf{desde hace} tres días en \textbf{Yaundé}.} OU \esp{\textbf{Hace} tres días \textbf{que} llueve en Yaundé.}
\textbf{3. « Depuis que » + fait réel.} \esp{\textbf{Desde que} \textbf{trabaja} en la radio, es famosa.} (Indicatif présent).
\textbf{4. Durée au passé (imparfait).} \esp{\textbf{Hacía} mucho tiempo \textbf{que} \textbf{esperaban} el autobús.}
\end{exoresolu}

\courssub{Synthèse comparative, pièges et entraînement thème/version}

\begin{methode}[Stratégie globale de thème (Français $\rightarrow$ Espagnol)]
\begin{enumerate}
\item \textbf{Lire la phrase entière} et \textbf{surligner} les trois notions cibles : « avoir beau », « demander », « depuis ».
\item \textbf{Traduire mentalement chaque segment} avec sa structure espagnole figée (ex: « il a beau » $\rightarrow$ \esp{por mucho que} + subj ; « il y a 2 mois que » $\rightarrow$ \esp{hace 2 meses que} + prés.).
\item \textbf{Assembler} en respectant l'ordre des mots espagnol (SVO, pronoms enclitiques/proclitiques, place de l'adjectif).
\item \textbf{Vérifier la check-list} : Accents (\esp{está, después, televisión, información, solución, Yaundé}), Accords (genre/nombre), Concordance (Subj. Présent vs Imparfait), Ponctuation (\esp{¿...? ¡...!}), Faux amis (\esp{demandar, realizar, asistir, éxito}).
\item \textbf{Relire à voix haute} : la phrase doit « sonner » espagnol.
\end{enumerate}
\end{methode}

\begin{exoresolu}[Thème complet type Bac -- Sujet : Médias et paix]
Traduisez : « Les journalistes ont beau poser des questions, le ministre ne répond pas. Il y a deux mois que les habitants demandent une nouvelle école, et depuis que la radio locale parle du problème, tout le village réclame une solution. »
\tcblower
\textbf{Découpage notionnel :}
\begin{itemize}
\item « ont beau poser » $\rightarrow$ \esp{Por mucho que} + \textbf{subjonctif présent} (\esp{hagan} / \esp{formulen}).
\item « posent des questions » $\rightarrow$ \esp{hacer preguntas} / \esp{formular preguntas} (éviter \esp{preguntar preguntas}).
\item « il y a deux mois que » $\rightarrow$ \esp{Hace dos meses que} + \textbf{présent}.
\item « demandent une école » (chose) $\rightarrow$ \esp{piden}.
\item « depuis que la radio parle » $\rightarrow$ \esp{Desde que la radio \textbf{habla}} (indicatif) + \textbf{présent}.
\item « réclame » $\rightarrow$ \esp{reclama} / \esp{exige} (éviter \esp{*demanda}).
\end{itemize}

\textbf{Traduction rédigée :}
\esp{Por mucho que los periodistas \textbf{hagan} preguntas, el ministro no \textbf{responde}. \textbf{Hace dos meses que} los habitantes \textbf{piden} una nueva escuela, y \textbf{desde que} la radio local \textbf{habla} del problema, todo el pueblo \textbf{reclama} una solución.}

\textbf{Justifications grammaticales :}
\begin{itemize}
\item \esp{hagan} : Subj. présent irrégulier de \esp{hacer} (obligatoire après \esp{por mucho que} + principale présent).
\item \esp{Hace dos meses que piden} : Présent espagnol pour durée continue (pas de \esp{*han pedido}).
\item \esp{desde que habla} : Indicatif (fait réel, constatable).
\item \esp{reclama} : Verbe précis pour « réclamer/exiger », évite le faux ami \esp{demandar}.
\item Accents : \esp{solución, radio, Yaundé} (pas d'accent sur \esp{radio}, accent sur \esp{Yaundé} usage hispanophone).
\end{itemize}
\end{exoresolu}

\begin{methode}[Stratégie globale de version (Espagnol $\rightarrow$ Français)]
\begin{enumerate}
\item \textbf{Repérer les marqueurs grammaticaux} : \esp{por mucho/más que}, \esp{hace/hacía... que}, \esp{desde hace}, \esp{desde que}, \esp{pedir/preguntar/rogar/solicitar}, modes/temps.
\item \textbf{Traduire chaque marqueur} par son équivalent français \textbf{idiomatique} (pas mot-à-mot) :
\begin{itemize}
\item \esp{Por mucho que} + subj. $\rightarrow$ « Il a beau / Quel que soit... »
\item \esp{Hace... que} + présent $\rightarrow$ « Cela fait ... que / Il y a ... que » + présent.
\item \esp{Hacía... que} + imparfait $\rightarrow$ « Cela faisait ... que / Il y avait ... que » + imparfait.
\item \esp{Pedir} $\rightarrow$ « demander (qch) » ; \esp{Preguntar} $\rightarrow$ « demander (info) / poser une question » ; \esp{Rogar} $\rightarrow$ « supplier / prier » ; \esp{Solicitar} $\rightarrow$ « solliciter / demander officiellement ».
\end{itemize}
\item \textbf{Rendre les temps/modes} : Subj. espagnol $\rightarrow$ Subj. français ; Imparfait $\rightarrow$ Imparfait ; Passé simple $\rightarrow$ Passé simple / Passé composé ; Plus-que-parfait pour antériorité.
\item \textbf{Produire un français naturel, élégant, sans calque.}
\end{enumerate}
\end{methode}

\begin{exoresolu}[Version type Bac -- Sujet : Désinformation et crise]
Traduisez en français :
\esp{« Por más que el gobierno \textbf{rogó} a los medios que \textbf{informaran} con calma, \textbf{hacía semanas que} las noticias falsas \textbf{circulaban} por las redes sociales. \textbf{Desde que} \textbf{estalló} la crisis, los ciudadanos \textbf{pedían} explicaciones claras. »}
\tcblower
\textbf{Repérage technique :}
\begin{itemize}
\item \esp{Por más que... rogó} : Concession passée (Passé simple \esp{rogó}) $\rightarrow$ « Le gouvernement eut beau... »
\item \esp{rogó que informaran} : \esp{rogar que} + \textbf{Subj. Imparfait} (\esp{informaran}) $\rightarrow$ « pria/demanda instamment que » + subj. français (\esp{informassent}) ou infinitif (\esp{d'informer}).
\item \esp{hacía semanas que circulaban} : Durée passée + Imparfait $\rightarrow$ « Il y avait des semaines que circulaient / Cela faisait des semaines que circulaient. »
\item \esp{Desde que estalló} : \esp{Desde que} + \textbf{Passé simple} (\esp{estalló}) $\rightarrow$ Antériorité par rapport à l'imparfait principal $\rightarrow$ \textbf{Plus-que-parfait} : « Depuis que la crise \textbf{eut éclaté / avait éclaté} ».
\item \esp{pedían} : Imparfait (description passé / action continue) $\rightarrow$ « demandaient ».
\end{itemize}

\textbf{Traduction rédigée :}
« Le gouvernement \textbf{eut beau prier} les médias d'\textbf{informer} avec calme, \textbf{il y avait des semaines que} les fausses nouvelles \textbf{circulaient} sur les réseaux sociaux. \textbf{Depuis que la crise avait éclaté}, les citoyens \textbf{demandaient} des explications claires. »

\textbf{Points de vigilance :}
\begin{itemize}
\item « Eut beau » (passé simple littéraire) ou « avait beau » (plus courant) rend la concession passée.
\item \esp{informaran} $\rightarrow$ « d'informer » (infinitif après « prier qqn de ») ou « qu'ils informassent » (subj. littéraire). L'infinitif est plus naturel.
\item \esp{noticias falsas} = « fausses nouvelles » ou « infox » (néologisme officiel) / « fake news » (anglicisme toléré). Pas « nouvelles fausses » (ordre adjectif).
\end{itemize}
\end{exoresolu}

\begin{attention}[Les 5 erreurs fatales au Bac -- Récapitulatif]
\begin{enumerate}
\item \textbf{\esp{Tener beau} / \esp{Tener bello}} : \textbf{Interdit}. $\rightarrow$ \esp{Por mucho que} / \esp{Por más que} + \textbf{Subjonctif} / \esp{Aunque}.
\item \textbf{Oublier le subjonctif} après \esp{Por mucho que / Por más que / Por + adj + que}. \esp{*Por mucho que estudia} $\rightarrow$ \esp{Por mucho que \textbf{estudie}}. Au passé : \esp{estudiara}.
\item \textbf{Confondre \esp{Pedir} / \esp{Preguntar} / \esp{Demandar}}.
\begin{itemize}
\item Je demande un livre : \esp{Pido un libro}.
\item Je demande l'heure : \esp{Pregunto la hora}.
\item Je poursuis en justice : \esp{Demando al conductor}.
\end{itemize}
\item \textbf{Passé composé avec « depuis » (action continue)}. \esp{*He vivido aquí desde hace 5 años} $\rightarrow$ \esp{\textbf{Vivo} aquí desde hace 5 años / \textbf{Hace} 5 años \textbf{que vivo} aquí}.
\item \textbf{Négliger accents et ponctuation} : \esp{está} (verbe) vs \esp{esta} (dét.) ; \esp{después, televisión, información, solución, Yaundé, también, país}. Questions : \esp{¿...?} ; Exclamations : \esp{¡...!}.
\end{enumerate}
\end{attention}

\begin{savaistu}[L'espagnol au Cameroun et en Guinée équatoriale]
La Guinée équatoriale (Malabo, Bata) est le seul pays africain où l'espagnol est langue officielle (avec le français et le portugais). L'espagnol y est la langue de l'administration, de l'éducation et des médias (\esp{TVGE, Radio Nacional}). Au Cameroun, l'espagnol est la LV2 la plus enseignée après l'anglais. De nombreux Camerounais font des études en Espagne ou en Amérique latine (Bourses \esp{MAEC-AECID}, \esp{Carolina}). Le \esp{español ecuatoguineen} a des particularités phonétiques (yeísmo fort, aspiration de la \esp{s} finale, usage de \esp{ustedes} comme unique 2e personne pluriel) et lexicales (emprunts aux langues bantu : \esp{nganga}, \esp{mbó}, \esp{akwá}). Connaître cette variété enrichit la compétence culturelle du bac.
\end{savaistu}

\begin{exercice}[Entraînement autonome -- Thème / Version]
\textbf{A. Thème (Français $\rightarrow$ Espagnol)}
1. Nous avons beau expliquer, ils ne comprennent pas.
2. Depuis quand habites-tu à Douala ? -- J'y habite depuis 2010.
3. Le délégué a demandé (officiellement) une audience au ministre.
4. Elle a beau être fatiguée, elle continue à travailler.
5. Cela fait une heure que nous attendons le prof.

\textbf{B. Version (Espagnol $\rightarrow$ Français)}
1. \esp{Por más que lo intentaron, no pudieron abrir la puerta.}
2. \esp{Desde que llegó el nuevo director, todo ha cambiado.}
3. \esp{Los estudiantes pidieron que se redujera la matrícula.}
4. \esp{Hace tres años que no voy a la playa.}
5. \esp{Te ruego que me digas la verdad.}
\end{exercice}

\begin{corrige}[Exercice E25]
\textbf{A. Thème}
1. \esp{Por mucho que \textbf{expliquemos}, no \textbf{entienden}.} (Subj. présent \esp{expliquemos} -- principale présent).
2. \esp{¿Desde cuándo vives en Duala? -- Vivo allí \textbf{desde} 2010.} (Point de départ).
3. \esp{El delegado \textbf{solicitó} una audiencia al ministro.} (Officiel).
4. \esp{Por \textbf{cansada} que \textbf{esté}, sigue \textbf{trabajando}.} (\esp{Por} + adj fém + \esp{que} + subj. \esp{estar} ; gérondif \esp{trabajando}).
5. \esp{\textbf{Hace} una hora \textbf{que} \textbf{esperamos} al profe.} (Durée + présent).

\textbf{B. Version}
1. « Ils eurent beau essayer / Quel que fut leur effort, ils ne purent ouvrir la porte. » (\esp{intentaron} passé simple $\rightarrow$ subj. imparfait implicite dans la concession passée ; \esp{pudieron} passé simple).
2. « Depuis que le nouveau directeur est arrivé, tout a changé. » (\esp{llegó} passé simple $\rightarrow$ passé composé français ; \esp{ha cambiado} parfait composé $\rightarrow$ passé composé).
3. « Les étudiants demandèrent que l'inscription soit réduite. » (\esp{pidieron que se redujera} $\rightarrow$ subj. imparfait \esp{redujera} $\rightarrow$ subj. présent \esp{soit réduite} ou imparfait \esp{fût réduite}).
4. « Cela fait trois ans que je ne vais pas à la plage / Je n vais pas à la plage depuis trois ans. » (\esp{Hace tres años que no voy} $\rightarrow$ présent + négation).
5. « Je te prie / Te supplie de me dire la vérité. » (\esp{Te ruego que me digas} $\rightarrow$ infinitif naturel).
\end{corrige}

\newpage

\coursec{Exercices}

\courssub{Je vérifie mes connaissances}

\begin{exercice}[QCM : la bonne traduction]
Pour chaque phrase française, choisis la bonne traduction espagnole parmi les quatre options proposées. Une seule réponse est correcte.

\begin{enumerate}
\item « Il a beau crier, personne ne l'entend. »
\begin{enumerate}
\item \esp{Por más que grita, nadie le oye.}
\item \esp{Por más que grite, nadie le oye.}
\item \esp{Por mucho gritar, nadie le oye.}
\item \esp{Aunque gritaba, nadie le oye.}
\end{enumerate}
\item « Elle demande un renseignement. »
\begin{enumerate}
\item \esp{Pide información.}
\item \esp{Pregunta una información.}
\item \esp{Pide por una información.}
\item \esp{Demanda una información.}
\end{enumerate}
\item « J'attends depuis une heure. »
\begin{enumerate}
\item \esp{Espero desde una hora.}
\item \esp{Espero desde hace una hora.}
\item \esp{Hace una hora espero.}
\item \esp{Espero desde que una hora.}
\end{enumerate}
\item « Depuis qu'il est parti, elle pleure. »
\begin{enumerate}
\item \esp{Desde hace que se fue, llora.}
\item \esp{Desde que se fue, llora.}
\item \esp{Hace que se fue, llora.}
\item \esp{Desde se fue, llora.}
\end{enumerate}
\end{enumerate}
\end{exercice}

\begin{exercice}[Vrai ou faux justifié]
Dis si les affirmations suivantes sont vraies ou fausses, puis justifie ta réponse en une phrase.

\begin{enumerate}
\item Après \esp{por más que} et \esp{por mucho que}, on emploie toujours le subjonctif, quel que soit le sens.
\item \esp{Preguntar} signifie « demander quelque chose » (un objet, une faveur).
\item \esp{Desde} seul s'emploie devant un point de départ précis (date, moment), tandis que \esp{desde hace} s'emploie devant une durée.
\item \esp{Rogar} est un verbe régulier qui se conjugue comme \esp{hablar}.
\item \esp{Hace diez años que vivo en Yaoundé} est une traduction correcte de « J'habite à Yaoundé depuis dix ans ».
\end{enumerate}
\end{exercice}

\begin{exercice}[Vocabulaire : les nuances de « demander »]
Complète chaque phrase avec le verbe qui convient, choisi parmi : \esp{pedir}, \esp{preguntar}, \esp{rogar}, \esp{solicitar}. Conjugue le verbe à la personne et au temps qui conviennent.

\begin{enumerate}
\item Los alumnos \dots\ al profesor que repita la explicación.
\item El periodista \dots\ al ministro cuándo visitará Douala.
\item Te \dots\ que me escuches un momento, por favor.
\item Los candidatos \dots\ un empleo en la empresa de cacao.
\item Mi hermana \dots\ perdón por llegar tarde.
\item El viajero \dots\ dónde está la estación de autobuses.
\end{enumerate}
\end{exercice}

\begin{exercice}[Conjugaison et modes à compléter]
Complète les phrases suivantes en mettant le verbe entre parenthèses à la forme qui convient (indicatif ou subjonctif, temps approprié).

\begin{enumerate}
\item Por mucho que \dots\ (estudiar, tú), no aprobarás sin trabajar.
\item Aunque \dots\ (llover) mucho ayer, el partido se jugó.
\item Por más que le \dots\ (rogar, yo), no quiso venir.
\item Te \dots\ (pedir, yo) un favor: préstame tu diccionario.
\item \dots\ (Hacer) dos semanas que esperamos una respuesta.
\item El niño \dots\ (preguntar) por qué el cielo es azul.
\end{enumerate}
\end{exercice}

\courssub{Je m'entraîne}

\begin{exercice}[Thème : phrases isolées]
Traduis en espagnol les phrases suivantes, en choisissant à chaque fois la structure la plus naturelle.

\begin{enumerate}
\item Il a beau crier, personne ne l'entend.
\item Elle demande un renseignement à l'accueil.
\item J'attends depuis une heure devant le lycée.
\item Depuis qu'il est parti, elle pleure tous les soirs.
\item Nous demandons un visa pour l'Espagne.
\item Il a beau être riche, il n'est pas heureux.
\end{enumerate}
\end{exercice}

\begin{exercice}[Transformation : depuis une date / depuis une durée]
Réécris chaque phrase en transformant l'expression temporelle selon le modèle indiqué entre parenthèses, sans changer le sens.

Modèle : \esp{Vivo aquí desde 2010.} $\rightarrow$ \esp{Vivo aquí desde hace quince años.} / \esp{Hace quince años que vivo aquí.}

\begin{enumerate}
\item Estudio español desde 2020. (transformer avec \esp{desde hace})
\item Trabaja en el hospital desde 2015. (transformer avec \esp{hace\ldots que})
\item Conozco a María desde hace cinco años. (transformer avec une date : 2020)
\item Hace dos meses que esperan el visado. (transformer avec \esp{desde hace})
\item No llueve desde abril. (transformer avec \esp{hace\ldots que} : six mois)
\end{enumerate}
\end{exercice}

\begin{exercice}[Texte à trous : les verbes de la demande]
Complète le texte suivant en mettant les verbes entre parenthèses à la forme et au mode qui conviennent.

\begin{textebox}[Un correo formal]
Estimado señor director:

Le escribo para \dots\ (solicitar) un puesto de prácticas en su emisora de radio. \dots\ (Hacer) dos semanas que envié mi primera carta y, por más que \dots\ (llamar, yo) a su secretaría, todavía no he recibido respuesta. Le \dots\ (rogar) que tenga en cuenta mi candidatura. Además, me gustaría \dots\ (preguntar) si es posible empezar en julio. Si usted \dots\ (poder) recibirme, se lo agradecería mucho. Espero que mi solicitud \dots\ (llegar) a buen puerto.

Atentamente,
André Mbarga, estudiante de terminale, Yaoundé.
\end{textebox}
\end{exercice}

\begin{exercice}[Appariements : structures et équivalents]
Associe chaque structure espagnole (colonne de gauche) à son équivalent français (colonne de droite). Attention, un intrus s'est glissé dans la colonne de droite.

\begin{center}
\begin{tabularx}{\linewidth}{|X|X|}
\hline
\rowcolor{popblueL} \textbf{Español} & \textbf{Français} \\
\hline
1. \esp{Por mucho que insista\ldots} & a. Il a beau insister\ldots \\
\hline
2. \esp{Aunque insiste\ldots} & b. Bien qu'il insiste (et c'est un fait)\ldots \\
\hline
3. \esp{Desde que llegó\ldots} & c. Depuis son arrivée\ldots \\
\hline
4. \esp{Desde su llegada\ldots} & d. Depuis qu'il est arrivé\ldots \\
\hline
5. \esp{Hace un año que llegó\ldots} & e. Il demande une explication. \\
\hline
 & f. Il est arrivé il y a un an. \\
\hline
\end{tabularx}
\end{center}
\end{exercice}

\courssub{Je me prépare au Bac}

\begin{exercice}[Compréhension et questions de langue]
Lis le texte suivant, puis réponds aux questions.

\begin{textebox}[La radio comunitaria de Bafut]
Desde hace diez años, la pequeña radio comunitaria de Bafut informa a los habitantes de la región. Por mucho que las autoridades han intentado modernizar los medios de comunicación, esta emisora sigue siendo la voz preferida de los campesinos. «Pedimos noticias sobre el precio del cacao y preguntamos por la salud de nuestros vecinos», explica Marta, la directora, que trabaja allí desde 2015. Por más que llueva o haga sol, los oyentes encienden sus aparatos cada mañana a las seis. Hace dos meses que la radio solicita una ayuda al ministerio, pero todavía no ha recibido respuesta. «Rogamos a los responsables que nos escuchen», concluye Marta, «porque desde que existe esta emisora, la gente se siente menos sola.»
\end{textebox}

\textbf{Questions de compréhension} (réponds en espagnol, avec des phrases complètes) :
\begin{enumerate}
\item ¿Cuánto tiempo hace que funciona la radio de Bafut?
\item ¿Qué piden los campesinos a la emisora?
\item ¿Qué ha solicitado la radio al ministerio? ¿Ha tenido éxito?
\item Según Marta, ¿qué efecto tiene la emisora sobre la gente?
\end{enumerate}

\textbf{Questions de langue} :
\begin{enumerate}
\item Relève dans le texte deux structures traduisant « avoir beau » et précise le mode du verbe qui suit.
\item Justifie l'emploi de \esp{desde hace diez años} et de \esp{desde 2015} dans le texte.
\item Pourquoi le subjonctif est-il employé après \esp{rogamos a los responsables que\ldots} ?
\end{enumerate}
\end{exercice}

\begin{exercice}[Thème et version type Baccalauréat]
\textbf{A. Thème.} Traduis en espagnol les phrases suivantes :

\begin{enumerate}
\item Il a beau répéter la question, l'élève ne comprend pas.
\item Les journalistes demandent une interview au ministre depuis trois semaines.
\item Bien qu'elle soit fatiguée, elle continue à travailler à la radio.
\item Je vous prie de bien vouloir m'envoyer le dossier.
\item Depuis que tu m'as écrit, je n'ai pas reçu d'autre nouvelle.
\end{enumerate}

\textbf{B. Version.} Traduis en français naturel le texte suivant (environ 60 mots) :

\begin{textebox}[Version]
Hace cinco años que Pablo dirige un periódico escolar en Malabo. Por mucho que pida ayuda a las empresas, nadie le responde. Desde que empezó el proyecto, los alumnos preguntan cada día por las noticias del barrio. «Solicitamos una subvención desde hace meses», explica Pablo, «y rogamos a todos que crean en los jóvenes.»
\end{textebox}
\end{exercice}

\begin{exercice}[Expression écrite : un correo formal]
Tu es élève en Terminale A à Douala. Tu as envoyé, il y a deux semaines, une demande de stage à une ONG espagnole qui travaille pour la culture de la paix. Malgré tes relances, tu n'as reçu aucune réponse.

Rédige en espagnol un courriel formel de \textbf{120 à 150 mots} dans lequel tu :
\begin{itemize}
\item sollicites officiellement un stage d'observation (utilise \esp{solicitar}) ;
\item rappelles que tu attends une réponse depuis deux semaines (utilise \esp{desde hace} ou \esp{hace\ldots que}) ;
\item expliques que tu as beau relancer, personne ne répond (utilise \esp{por más que} ou \esp{por mucho que} + subjonctif) ;
\item pries poliment le directeur de te répondre rapidement (utilise \esp{rogar}) ;
\item poses une question sur la date de début du stage (utilise \esp{preguntar}).
\end{itemize}

Respecte la structure d'un courriel formel : formule d'appel, introduction, développement, formule de politesse finale, signature.
\end{exercice}

\begin{methode}[Réussir le courriel formel]
\begin{enumerate}
\item \textbf{Appel} : \esp{Estimado señor director:} / \esp{Estimada señora directora:} (deux-points, puis majuscule à la ligne suivante).
\item \textbf{Introduction} : annonce l'objet avec \esp{solicitar} : \esp{Le escribo para solicitar un puesto de prácticas de observación en su organización.}
\item \textbf{Rappel temporel} : \esp{Hace dos semanas que envié mi solicitud} ou \esp{Espero una respuesta desde hace dos semanas.}
\item \textbf{Concession} : \esp{Por más que he llamado y he escrito, nadie me ha respondido.} (ou subjonctif : \esp{Por más que llame\ldots})
\item \textbf{Demande polie} : \esp{Le ruego que me responda lo antes posible.} (subjonctif obligatoire après \esp{rogar que}).
\item \textbf{Question} : \esp{Quisiera preguntar cuándo podría empezar el período de prácticas.}
\item \textbf{Formule finale} : \esp{A la espera de su respuesta, le saluda atentamente,} puis prénom, nom, classe, ville.
\end{enumerate}
\end{methode}

\newpage

\coursec{Corrigés des exercices}

\begin{corrige}[Exercice 1 — QCM : la bonne traduction]
\begin{enumerate}
\item \textbf{b.} \esp{Por más que grite, nadie le oye.} Après \esp{por más que} / \esp{por mucho que} exprimant la concession, le subjonctif présent est la forme standard quand le fait n'est pas présenté comme réel. L'option a (indicatif) n'est pas la réponse attendue ; l'option c est agrammaticale (\esp{por mucho} ne se construit pas avec l'infinitif) ; l'option d emploie un imparfait inadapté.
\item \textbf{a.} \esp{Pide información.} « Demander quelque chose » = \esp{pedir} + complément direct, sans préposition. \esp{Preguntar} = poser une question ; \esp{demandar} est un faux ami (exiger, poursuivre en justice).
\item \textbf{b.} \esp{Espero desde hace una hora.} Devant une durée, on emploie \esp{desde hace} ; \esp{desde} seul introduit un point de départ.
\item \textbf{b.} \esp{Desde que se fue, llora.} Devant une proposition verbale, on emploie \esp{desde que} + indicatif.
\end{enumerate}
\end{corrige}

\begin{corrige}[Exercice 2 — Vrai ou faux justifié]
\begin{enumerate}
\item \textbf{Faux.} Quand le fait est présenté comme réel, l'indicatif est possible : \esp{Por mucho que insiste, no lo consigue.} Le subjonctif domine quand la concession est hypothétique ou générale.
\item \textbf{Faux.} \esp{Preguntar} signifie « poser une question » ; « demander quelque chose » se dit \esp{pedir}.
\item \textbf{Vrai.} \esp{desde 2015}, \esp{desde abril} (point de départ) ; \esp{desde hace diez años} (durée écoulée jusqu'à aujourd'hui).
\item \textbf{Faux.} \esp{Rogar} est irrégulier : alternance \esp{o} $\rightarrow$ \esp{ue} (\esp{ruego, ruegas, ruega, ruegan}) et \esp{o} $\rightarrow$ \esp{u} au subjonctif présent pour \esp{nosotros/vosotros} (\esp{roguemos, roguéis}).
\item \textbf{Vrai.} La structure \esp{hace + durée + que + présent} équivaut exactement à « depuis + durée » : \esp{Hace diez años que vivo en Yaoundé.}
\end{enumerate}
\end{corrige}

\begin{corrige}[Exercice 3 — Vocabulaire : les nuances de « demander »]
\begin{enumerate}
\item \esp{piden} (\esp{pedir algo a alguien} ; noter le subjonctif \esp{repita} après \esp{pedir que}).
\item \esp{pregunta} (poser une question : \esp{preguntar cuándo\ldots}).
\item \esp{ruego} (prier avec insistance : \esp{rogar que} + subjonctif ; alternance \esp{o} $\rightarrow$ \esp{ue}).
\item \esp{solicitan} (demande officielle, administrative : un emploi, un visa, une bourse).
\item \esp{pide} (\esp{pedir perdón}, expression figée).
\item \esp{pregunta} (s'informer : \esp{preguntar dónde\ldots}).
\end{enumerate}
\end{corrige}

\begin{corrige}[Exercice 4 — Conjugaison et modes à compléter]
\begin{enumerate}
\item \esp{estudies} : subjonctif présent après \esp{por mucho que} (concession hypothétique).
\item \esp{llovió} : indicatif passé simple après \esp{aunque}, car le fait est réel et daté (\esp{ayer}).
\item \esp{rogué} ou \esp{rogara} : l'indicatif passé simple présente l'insistance comme un fait accompli ; l'imparfait du subjonctif insiste sur l'inutilité de la démarche. Les deux sont corrects.
\item \esp{pido} : présent de \esp{pedir} (alternance \esp{e} $\rightarrow$ \esp{i}).
\item \esp{Hace} : \esp{hacer} impersonnel à la 3\up{e} personne du singulier dans la structure \esp{hace + durée + que}.
\item \esp{pregunta} : présent de vérité générale ; noter l'accent dans la question indirecte \esp{por qué}.
\end{enumerate}
\end{corrige}

\begin{corrige}[Exercice 5 — Thème : phrases isolées]
\begin{enumerate}
\item \esp{Por más que grite, nadie le oye.} (ou \esp{Por mucho que grite\ldots})
\item \esp{Pide información en recepción.} (ou \esp{Pide un dato en la oficina de información.})
\item \esp{Espero desde hace una hora delante del liceo.} (ou \esp{Hace una hora que espero delante del liceo.})
\item \esp{Desde que se fue, llora todas las noches.}
\item \esp{Solicitamos un visado para España.} (demande officielle $\rightarrow$ \esp{solicitar})
\item \esp{Por mucho que sea rico, no es feliz.} (ou \esp{Aunque es rico, no es feliz.} si l'on constate le fait)
\end{enumerate}
\textbf{Points de vigilance} : ne jamais écrire \esp{pedir por} ; ne pas oublier le subjonctif après \esp{por más que} ; bien distinguer durée (\esp{desde hace}) et point de départ (\esp{desde}).
\end{corrige}

\begin{corrige}[Exercice 6 — Transformation : depuis une date / depuis une durée]
\begin{enumerate}
\item \esp{Estudio español desde hace cinco años.}
\item \esp{Hace diez años que trabaja en el hospital.}
\item \esp{Conozco a María desde 2020.}
\item \esp{Esperan el visado desde hace dos meses.}
\item \esp{Hace seis meses que no llueve.} (attention : la négation doit être conservée)
\end{enumerate}
\textbf{Points de vigilance} : dans la structure \esp{hace\ldots que}, le verbe reste au présent ; \esp{desde hace} se place après le verbe ou en tête de phrase, jamais entre \esp{desde} et le nom de durée.
\end{corrige}

\begin{corrige}[Exercice 7 — Texte à trous : les verbes de la demande]
\esp{solicitar} (infinitif après \esp{para}) ; \esp{Hace} (impersonnel) ; \esp{llame} (subjonctif présent après \esp{por más que} : j'ai beau appeler) ; \esp{ruego} (présent, alternance \esp{o} $\rightarrow$ \esp{ue} ; \esp{rogar que} + subjonctif \esp{tenga}) ; \esp{preguntar} (infinitif après \esp{me gustaría}) ; \esp{puede} (indicatif après \esp{si} conditionnel) ; \esp{llegue} (subjonctif après \esp{esperar que}).

\textbf{Texte complété} : \esp{Le escribo para solicitar un puesto de prácticas en su emisora de radio. Hace dos semanas que envié mi primera carta y, por más que llame a su secretaría, todavía no he recibido respuesta. Le ruego que tenga en cuenta mi candidatura. Además, me gustaría preguntar si es posible empezar en julio. Si usted puede recibirme, se lo agradecería mucho. Espero que mi solicitud llegue a buen puerto.}
\end{corrige}

\begin{corrige}[Exercice 8 — Appariements : structures et équivalents]
1 -- a ; 2 -- b ; 3 -- d ; 4 -- c ; 5 -- f. L'intrus est \textbf{e} : « Il demande une explication » se traduirait \esp{Pide una explicación}, structure absente de la colonne de gauche.

\textbf{Points de vigilance} : bien distinguer \esp{desde que} + verbe conjugué (proposition : « depuis qu'il est arrivé ») et \esp{desde} + nom (\esp{desde su llegada} : « depuis son arrivée ») ; \esp{hace un año que llegó} avec un passé simple signifie « il est arrivé il y a un an », et non « depuis un an ».
\end{corrige}

\begin{corrige}[Exercice 9 — Compréhension et questions de langue]
\textbf{Compréhension} (réponses attendues) :
\begin{enumerate}
\item \esp{Hace diez años que funciona la radio de Bafut.}
\item \esp{Los campesinos piden noticias sobre el precio del cacao y preguntan por la salud de sus vecinos.}
\item \esp{La radio ha solicitado una ayuda al ministerio, pero no ha tenido éxito: todavía no ha recibido respuesta.}
\item \esp{Según Marta, gracias a la emisora la gente se siente menos sola.}
\end{enumerate}

\textbf{Langue} :
\begin{enumerate}
\item \esp{Por mucho que las autoridades han intentado\ldots} : indicatif (parfait composé), car les tentatives sont présentées comme un fait réel ; \esp{Por más que llueva o haga sol\ldots} : subjonctif présent, car il s'agit d'une éventualité générale (« qu'il pleuve ou qu'il fasse beau »).
\item \esp{Desde hace diez años} introduit une durée (dix ans écoulés jusqu'à aujourd'hui) ; \esp{desde 2015} introduit un point de départ précis (une date).
\item \esp{Rogar que} exprime une demande portant sur l'action d'autrui : comme tous les verbes de volonté et de demande (\esp{pedir que, querer que, desear que}), il exige le subjonctif dans la subordonnée (\esp{escuchen}).
\end{enumerate}
\textbf{Barème indicatif} : compréhension 4 $\times$ 1 pt ; questions de langue 3 $\times$ 2 pts (1 pt pour la réponse, 1 pt pour la justification).
\end{corrige}

\begin{corrige}[Exercice 10 — Thème et version type Baccalauréat]
\textbf{A. Thème} (propositions de corrigé) :
\begin{enumerate}
\item \esp{Por más que repita la pregunta, el alumno no comprende.}
\item \esp{Los periodistas piden una entrevista al ministro desde hace tres semanas.} (ou \esp{Hace tres semanas que los periodistas piden\ldots})
\item \esp{Aunque está cansada, sigue trabajando en la radio.} (fait réel $\rightarrow$ indicatif ; \esp{Aunque esté cansada} est possible si la fatigue est hypothétique)
\item \esp{Le ruego que me envíe el expediente.} (ou \esp{Les ruego que me envíen el dossier.})
\item \esp{Desde que me escribiste, no he recibido ninguna otra noticia.}
\end{enumerate}

\textbf{B. Version} (proposition de traduction) :

« Cela fait cinq ans que Pablo dirige un journal scolaire à Malabo. Il a beau demander de l'aide aux entreprises, personne ne lui répond. Depuis que le projet a commencé, les élèves s'informent chaque jour des nouvelles du quartier. "Nous sollicitons une subvention depuis des mois, explique Pablo, et nous prions tout le monde de croire en la jeunesse." »

\textbf{Barème indicatif} : thème 5 $\times$ 2 pts (1 pt pour la structure ciblée, 1 pt pour l'exactitude globale) ; version 10 pts (fidélité au sens 5 pts, qualité du français 3 pts, rendu des trois structures de la leçon 2 pts).

\textbf{Points de vigilance} : « répéter » = \esp{repetir} (alternance \esp{e} $\rightarrow$ \esp{i} : \esp{repita}) ; « je vous prie de\ldots » se rend élégamment par \esp{le ruego que} + subjonctif ; en version, \esp{por mucho que pida} doit être rendu par « avoir beau » et non par « bien que » si l'on veut conserver l'insistance.
\end{corrige}

\begin{corrige}[Exercice 11 — Expression écrite : un correo formal]
\textbf{Proposition de rédaction modèle} (environ 135 mots) :

\esp{Estimado señor director:}

\esp{Le escribo para solicitar un puesto de prácticas de observación en su organización, dedicada a la cultura de paz. Hace dos semanas que envié mi primera solicitud y espero una respuesta desde hace quince días. Por más que he llamado a su oficina y he enviado varios correos, nadie me ha respondido todavía. Le ruego que tenga en cuenta mi candidatura y que me responda lo antes posible. Además, quisiera preguntar cuándo podría empezar el período de prácticas, porque debo organizar mi viaje desde Douala. Soy alumno de Terminale A y me apasionan los medios de comunicación. A la espera de su respuesta, le saluda atentamente,}

\esp{Jean Fotso, alumno de Terminale A, Douala (Camerún).}

\textbf{Barème indicatif} (sur 20) : respect des cinq consignes et des structures imposées 5 pts ; correction grammaticale (subjonctif après \esp{rogar que} et \esp{por más que}, \esp{desde hace} bien employé) 6 pts ; richesse du lexique 3 pts ; structure du courriel formel (appel, introduction, formule finale, signature) 4 pts ; nombre de mots respecté et présentation 2 pts.

\textbf{Points de vigilance} : après la formule d'appel, deux-points puis majuscule à la ligne ; \esp{rogar} exige le subjonctif (\esp{que me responda}) ; ne pas écrire \esp{desde dos semanas} pour une durée mais \esp{desde hace dos semanas} ; éviter le calque \esp{demandar un stage} et préférer \esp{solicitar unas prácticas}.
\end{corrige}

\newpage

\coursec{Fiche bilan}

\begin{popfigure}
\begin{tikzpicture}[
  mindmap,
  grow cyclic,
  every node/.style={concept, align=center, font=\small},
  root concept/.append style={concept color=popblue, font=\small\bfseries, text=white, minimum size=2.6cm},
  level 1 concept/.append style={sibling angle=51, level distance=3.6cm, minimum size=2.2cm, font=\scriptsize\bfseries, text=white},
  level 2 concept/.append style={sibling angle=45, level distance=2.4cm, minimum size=1.7cm, font=\scriptsize, text=popdark},
]
\node[root concept, align=center]{E25\\Traduction\\avoir beau /\\demander / depuis}
  child[concept color=poporange]{ node[align=center]{Avoir beau :\\concession}
    child[concept color=poporangeL]{ node[align=center]{por más que\\+ subj.} }
    child[concept color=poporangeL]{ node[align=center]{por mucho\\que + subj.} }
    child[concept color=poporangeL]{ node[align=center]{aunque\\+ ind./subj.} }
  }
  child[concept color=popgreen]{ node[align=center]{Demander :\\4 verbes}
    child[concept color=popgreenL]{ node[align=center]{pedir\\(réclamer)} }
    child[concept color=popgreenL]{ node[align=center]{preguntar\\(interroger)} }
    child[concept color=popgreenL]{ node[align=center]{rogar\\(prier)} }
    child[concept color=popgreenL]{ node[align=center]{solicitar\\(formel)} }
  }
  child[concept color=poppurple]{ node[align=center]{Depuis :\\temps}
    child[concept color=poppurpleL]{ node[align=center]{desde\\(point de départ)} }
    child[concept color=poppurpleL]{ node[align=center]{desde hace\\+ durée} }
    child[concept color=poppurpleL]{ node[align=center]{hace... que\\+ présent} }
  }
  child[concept color=poppink]{ node[align=center]{Pièges\\francophones}
    child[concept color=poppinkL]{ node[align=center]{pedir $\neq$\\preguntar} }
    child[concept color=poppinkL]{ node[align=center]{depuis +\\présent =\\présent esp.} }
  }
  child[concept color=popgold]{ node[align=center]{Méthode\\thème/version}
    child[concept color=popgoldL]{ node[align=center]{Repérer le\\sens exact} }
    child[concept color=popgoldL]{ node[align=center]{Vérifier\\mode et temps} }
  };
\end{tikzpicture}
\legende{Carte mentale de la leçon E25 : trois notions clés de traduction.}
\end{popfigure}

\begin{bilan}
\textbf{Lexique clé à retenir}
\begin{itemize}
  \item \esp{pedir algo a alguien} : demander quelque chose à quelqu'un (\esp{Pedí ayuda a mi hermano.} « J'ai demandé de l'aide à mon frère. »)
  \item \esp{preguntar (por)} : demander, poser une question (\esp{Preguntó por el precio del cacao.} « Il a demandé le prix du cacao. »)
  \item \esp{rogar (u$\to$ue)} : prier, supplier (\esp{Te ruego que me escuches.} « Je te demande de m'écouter. »)
  \item \esp{solicitar} : solliciter, demander officiellement (\esp{Solicitó una beca.} « Il a demandé une bourse. »)
  \item \esp{por más / por mucho que} + subjonctif : avoir beau (\esp{Por más que estudie, no aprueba.} « Il a beau étudier, il échoue. »)
  \item \esp{desde} : depuis (point de départ) ; \esp{desde hace} : depuis (durée) ; \esp{hace... que} : cela fait... que
\end{itemize}

\textbf{Règles de grammaire à connaître par cœur}

\begin{center}
\begin{tabularx}{\linewidth}{|X|X|X|}
\hline
\rowcolor{popblueL} Français & Español & Remarque \\
\hline
Il a beau travailler... & \esp{Por mucho que trabaje...} & subjonctif obligatoire \\
\hline
Il a beau être riche... & \esp{Por muy rico que sea...} & \esp{por muy} + adjectif + \esp{que} \\
\hline
Bien qu'il pleuve... & \esp{Aunque llueva...} & subj. = hypothèse ; indic. = fait réel \\
\hline
Demander qqch à qqn & \esp{pedir algo a alguien} & jamais \esp{demandar} (faux ami) \\
\hline
Demander si / pourquoi & \esp{preguntar si / por qué} & question, information \\
\hline
Je vis ici depuis 2010. & \esp{Vivo aquí desde 2010.} & point de départ + présent \\
\hline
Je vis ici depuis dix ans. & \esp{Vivo aquí desde hace diez años.} & durée + présent \\
\hline
Cela fait dix ans que... & \esp{Hace diez años que vivo aquí.} & présent après \esp{que} \\
\hline
\end{tabularx}
\end{center}

\textbf{Méthodes express}
\begin{enumerate}
  \item « Avoir beau » : repérer la concession, choisir \esp{por mucho que} + subjonctif (solution la plus sûre au Bac).
  \item « Demander » : objet demandé $\to$ \esp{pedir} ; question posée $\to$ \esp{preguntar} ; supplication $\to$ \esp{rogar} ; contexte administratif $\to$ \esp{solicitar}.
  \item « Depuis » : date ou événement précis $\to$ \esp{desde} ; durée $\to$ \esp{desde hace} ou \esp{hace... que} ; toujours le présent espagnol si l'action continue.
\end{enumerate}
\end{bilan}

\begin{aretenir}[Checklist avant le Bac]
\begin{itemize}
  \item $\square$ Je sais traduire « avoir beau » par \esp{por mucho que} + subjonctif.
  \item $\square$ Je sais utiliser la structure \esp{por muy} + adjectif + \esp{que} + subjonctif.
  \item $\square$ Je distingue \esp{aunque} + indicatif (fait réel) et \esp{aunque} + subjonctif (hypothèse).
  \item $\square$ Je ne confonds jamais \esp{pedir} et \esp{preguntar}.
  \item $\square$ Je sais que \esp{demandar} est un faux ami (« poursuivre en justice »).
  \item $\square$ Je conjugue correctement \esp{rogar} (changement \esp{o $\to$ ue}).
  \item $\square$ J'emploie \esp{solicitar} uniquement dans un contexte formel ou administratif.
  \item $\square$ Je choisis \esp{desde} pour un point de départ et \esp{desde hace} pour une durée.
  \item $\square$ Je mets le verbe espagnol au présent après « depuis » quand l'action continue.
  \item $\square$ Je sais construire \esp{hace + durée + que} + présent.
  \item $\square$ Je vérifie toujours le mode (indicatif ou subjonctif) avant de rendre ma copie.
\end{itemize}
\end{aretenir}

\findecours

\end{document}
